% Generated by docs/cartographer/blueprint_gen.py. DO NOT EDIT.
% PX_RH_CLAIM: NOT_MADE
\section{Definitions and faithfulness interface}
\textbf{Route status: CHALLENGER / NOT\_RH.}\par
\textbf{PX\_RH\_CLAIM: NOT\_MADE.}\par
The project proposition \texttt{Q3.RH} is the classical open-strip Riemann Hypothesis over Mathlib's \texttt{riemannZeta}: every zero with real part strictly between zero and one has real part one half.
\begin{definition}[Project RH proposition]\label{def:q3-rh}
\lean{Q3.RH}
\texttt{Q3.RH} is defined in \texttt{Q3/Basic/Defs.lean}. This records definition identity and does not assert a proof.
\end{definition}
\begin{theorem}[Centered-Xi faithfulness bridge]\label{thm:rh-centered-xi}
\lean{Q3.RouteB.rh_iff_centeredXi_zeros_real}
\leanok
\mbox{}\par\smallskip
\begin{Verbatim}[breaklines=true,breakanywhere=true,fontsize=\small]
theorem rh_iff_centeredXi_zeros_real  : Q3.RH ↔ CenteredXiZerosReal
\end{Verbatim}
\smallskip
\end{theorem}
\begin{proof}\leanok The proof term is checked by Lean.\end{proof}
This equivalence of formulations is not a proof of either side.
\section{Conditional roof}
\begin{theorem}[Canonical strip-slot assembly]\label{thm:conditional-roof}
\lean{Q3.RouteB.CanonicalRHRoute.rh_of_canonical_strip_slots}
\leanok
\mbox{}\par\smallskip
\begin{Verbatim}[breaklines=true,breakanywhere=true,fontsize=\small]
theorem rh_of_canonical_strip_slots {Index : Type*} (C : CanonicalApproximation Index) (H2aAt S1At : Index → Prop) (anchor : ℂ) (hH1 : SlotH1 C) (hH2a : SlotH2a C H2aAt) (hanchor : SlotAnchor C anchor) (hS1 : SlotS1 C S1At) (hMontel : MontelAnchorGate C H2aAt S1At anchor) (h510 : Theorem510RealZeroBridge C H2aAt) (hS2 : SlotS2 C) : Q3.RH
\end{Verbatim}
\smallskip
\end{theorem}
\begin{proof}\leanok The proof term is checked by Lean.\end{proof}
The conclusion is conditional on every named hypothesis in the statement.

\section{Pillar G3 — de Branges/\allowbreak{}CvS port}

\begin{lemma}[G3\_\allowbreak{}CVS\_\allowbreak{}PORT / step 1]\label{assembly:g3-cvs-port-1-7b97a7626d07}
\notready
\textbf{Publication state:} \texttt{OPEN\_\allowbreak{}MATH}.\par
P1: GROUND\_\allowbreak{}CANONICAL\_\allowbreak{}PSTAR\_\allowbreak{}VALUE\_\allowbreak{}CROSSWALK — ground-значная CanonicalApproximation ЛИБО точное равенство selected-функции и (ненулевой множитель)×proposition59CCMTransform ground-строки\par
\textbf{Open receipt:} assembly status GAP.\par
\textbf{Recorded note:} вердикт батча 19.08: CvS подходит ПО ТИПУ, не по значению (trial против ground); постулировать запрещено\par
\end{lemma}

\begin{lemma}[G3\_\allowbreak{}CVS\_\allowbreak{}PORT / step 2]\label{assembly:g3-cvs-port-2-000e859e6268}
\notready
\textbf{Publication state:} \texttt{OPEN\_\allowbreak{}MATH}.\par
P2: Theorem510RealZeroBridge\_\allowbreak{}of\_\allowbreak{}groundP59 — сборка моста из CvS-движка\par
\textbf{Open receipt:} assembly status GAP.\par
\textbf{Recorded note:} сборка после P1; CvS-доказательство НЕ реформализуется\par
\end{lemma}

\section{Pillar G5 — critical-moment budget}

\begin{lemma}[G5\_\allowbreak{}CRITICAL\_\allowbreak{}MOMENT / step 1]\label{assembly:g5-critical-moment-1-906c85d62e13}
\notready
\textbf{Publication state:} \texttt{OPEN\_\allowbreak{}MATH}.\par
равномерный по k моментный бюджет: ∃C ∀k centeredCriticalMoment ≤ C·|rawFplus 0|\par
\textbf{Open receipt:} assembly status GAP.\par
\textbf{Recorded note:} мера ядром 19.08 (apply?): единственный аналитический канат G5; PairCofinal приезжает полем пакета D (шаг 25 GOAL057)\par
\end{lemma}

\section{Pillar G6 — continuum numerator and edge}

\begin{lemma}[GOAL057\_\allowbreak{}CONTINUUM\_\allowbreak{}NUMERATOR\_\allowbreak{}BRIDGE / step 0]\label{assembly:goal057-continuum-numerator-bridge-0-c940ca4a76d4}
\notready
\textbf{Publication state:} \texttt{READY\_\allowbreak{}WITHOUT\_\allowbreak{}EXACT\_\allowbreak{}DECLARATION\_\allowbreak{}RECEIPT}.\par
СТАРТ: проба не является источниковым объектом\par
\textbf{Assembly supplier field:} \texttt{аудит GOAL057\_\allowbreak{}ACTUAL\_\allowbreak{}NUMERATOR\_\allowbreak{}SOURCE\_\allowbreak{}TARGET\_\allowbreak{}AUDIT\_\allowbreak{}2026-08-07}.\par
\textbf{Open receipt:} no unique exact name in aristotle\_\allowbreak{}proofs.db.\par
\textbf{Recorded note:} вердикт PROBE\_\allowbreak{}NOT\_\allowbreak{}SOURCE\_\allowbreak{}TRIAL: 'numerically extremely close, but different constructions'. Классическая карта C10 — суррогат по форме. Это и запустило всю цепочку.\par
\end{lemma}

\begin{lemma}[GOAL057\_\allowbreak{}CONTINUUM\_\allowbreak{}NUMERATOR\_\allowbreak{}BRIDGE / step 1]\label{assembly:goal057-continuum-numerator-bridge-1-5f93b3b94e70}
\notready
\textbf{Publication state:} \texttt{READY\_\allowbreak{}WITHOUT\_\allowbreak{}EXACT\_\allowbreak{}DECLARATION\_\allowbreak{}RECEIPT}.\par
B1: привязка источникового ряда коэффициентов\par
\textbf{Assembly supplier field:} \texttt{D0PstarCCMFiniteSourceResidual}.\par
\textbf{Open receipt:} no unique exact name in aristotle\_\allowbreak{}proofs.db.\par
\textbf{Recorded note:} 6 теорем, 0 sorry. Шапка сама перечисляет НЕсделанное: не отождествляет конечное действие со сжатием континуального оператора Вейля, не даёт projection-leakage rate, не закрывает H4a1b.\par
\end{lemma}

\begin{lemma}[GOAL057\_\allowbreak{}CONTINUUM\_\allowbreak{}NUMERATOR\_\allowbreak{}BRIDGE / step 2]\label{assembly:goal057-continuum-numerator-bridge-2-4dbfee34f7ff}
\notready
\textbf{Publication state:} \texttt{READY\_\allowbreak{}WITHOUT\_\allowbreak{}EXACT\_\allowbreak{}DECLARATION\_\allowbreak{}RECEIPT}.\par
B2: перенос в операторный носитель (Riesz)\par
\textbf{Assembly supplier field:} \texttt{D0PstarCCMFiniteRieszOperator}.\par
\textbf{Open receipt:} no unique exact name in aristotle\_\allowbreak{}proofs.db.\par
\textbf{Recorded note:} 1 теорема, 0 sorry. НЕ определяет ambient A\_\allowbreak{}m, не доказывает членство в Dom(A\_\allowbreak{}m), не даёт ambient compression, не производит continuum numerator.\par
\end{lemma}

\begin{lemma}[GOAL057\_\allowbreak{}CONTINUUM\_\allowbreak{}NUMERATOR\_\allowbreak{}BRIDGE / step 3]\label{assembly:goal057-continuum-numerator-bridge-3-d39f0eee5b6e}
\notready
\textbf{Publication state:} \texttt{READY\_\allowbreak{}WITHOUT\_\allowbreak{}EXACT\_\allowbreak{}DECLARATION\_\allowbreak{}RECEIPT}.\par
B3.0A: формула Фурье для мод\par
\textbf{Assembly supplier field:} \texttt{D0PstarVModeFourierFormula}.\par
\textbf{Open receipt:} no unique exact name in aristotle\_\allowbreak{}proofs.db.\par
\end{lemma}

\begin{lemma}[GOAL057\_\allowbreak{}CONTINUUM\_\allowbreak{}NUMERATOR\_\allowbreak{}BRIDGE / step 4]\label{assembly:goal057-continuum-numerator-bridge-4-e33d197d338e}
\notready
\textbf{Publication state:} \texttt{READY\_\allowbreak{}WITHOUT\_\allowbreak{}EXACT\_\allowbreak{}DECLARATION\_\allowbreak{}RECEIPT}.\par
B3.0B1: взвешенная L2-норма мод\par
\textbf{Assembly supplier field:} \texttt{D0PstarVModeLogWeightedL2}.\par
\textbf{Open receipt:} no unique exact name in aristotle\_\allowbreak{}proofs.db.\par
\end{lemma}

\begin{lemma}[GOAL057\_\allowbreak{}CONTINUUM\_\allowbreak{}NUMERATOR\_\allowbreak{}BRIDGE / step 5]\label{assembly:goal057-continuum-numerator-bridge-5-557115991254}
\notready
\textbf{Publication state:} \texttt{READY\_\allowbreak{}WITHOUT\_\allowbreak{}EXACT\_\allowbreak{}DECLARATION\_\allowbreak{}RECEIPT}.\par
B3.0B2: доминирование точного арх. символа\par
\textbf{Assembly supplier field:} \texttt{D0PstarExactArchSymbolLogDomination}.\par
\textbf{Open receipt:} no unique exact name in aristotle\_\allowbreak{}proofs.db.\par
\end{lemma}

\begin{lemma}[GOAL057\_\allowbreak{}CONTINUUM\_\allowbreak{}NUMERATOR\_\allowbreak{}BRIDGE / step 6]\label{assembly:goal057-continuum-numerator-bridge-6-627fa4733339}
\notready
\textbf{Publication state:} \texttt{READY\_\allowbreak{}WITHOUT\_\allowbreak{}EXACT\_\allowbreak{}DECLARATION\_\allowbreak{}RECEIPT}.\par
B3.0B3: точный арх. символ в взвешенном L2\par
\textbf{Assembly supplier field:} \texttt{D0PstarExactArchSymbolWeightedModeL2}.\par
\textbf{Open receipt:} no unique exact name in aristotle\_\allowbreak{}proofs.db.\par
\end{lemma}

\begin{lemma}[GOAL057\_\allowbreak{}CONTINUUM\_\allowbreak{}NUMERATOR\_\allowbreak{}BRIDGE / step 7]\label{assembly:goal057-continuum-numerator-bridge-7-d5757ce20293}
\notready
\textbf{Publication state:} \texttt{READY\_\allowbreak{}WITHOUT\_\allowbreak{}EXACT\_\allowbreak{}DECLARATION\_\allowbreak{}RECEIPT}.\par
B3.0C: интегрируемость спаривания мод\par
\textbf{Assembly supplier field:} \texttt{D0PstarSourceArchModePairingIntegrable}.\par
\textbf{Open receipt:} no unique exact name in aristotle\_\allowbreak{}proofs.db.\par
\end{lemma}

\begin{lemma}[GOAL057\_\allowbreak{}CONTINUUM\_\allowbreak{}NUMERATOR\_\allowbreak{}BRIDGE / step 8]\label{assembly:goal057-continuum-numerator-bridge-8-944251764855}
\notready
\textbf{Publication state:} \texttt{READY\_\allowbreak{}WITHOUT\_\allowbreak{}EXACT\_\allowbreak{}DECLARATION\_\allowbreak{}RECEIPT}.\par
B3.0D: эрмитовость ядра спаривания\par
\textbf{Assembly supplier field:} \texttt{D0PstarSourceArchModePairingKernel}.\par
\textbf{Open receipt:} no unique exact name in aristotle\_\allowbreak{}proofs.db.\par
\textbf{Recorded note:} только арх. ядро фиксированных мод; НЕ отождествляет значение интеграла, элемент CCM-матрицы, полную источниковую форму Вейля или граф оператора.\par
\end{lemma}

\begin{lemma}[GOAL057\_\allowbreak{}CONTINUUM\_\allowbreak{}NUMERATOR\_\allowbreak{}BRIDGE / step 9]\label{assembly:goal057-continuum-numerator-bridge-9-4a396ec287ae}
\notready
\textbf{Publication state:} \texttt{READY\_\allowbreak{}WITHOUT\_\allowbreak{}EXACT\_\allowbreak{}DECLARATION\_\allowbreak{}RECEIPT}.\par
B3.0E: crosswalk CCM <-> WR (E1,E2,E3,E4A,E4B1,E4B2,E4C)\par
\textbf{Assembly supplier field:} \texttt{закрыто 08.08 16:47 'Close all-mode CCM-WR crosswalk'}.\par
\textbf{Open receipt:} READY has no exact supplier name and source path.\par
\textbf{Recorded note:} разветвился на 7 подшагов за день. E1 скалярное тождество, E2 Fubini, E3 корреляция, E4A внедиагональ, E4B1 endpoint ledger, E4B2 диагональ, E4C все моды. | WR-ЧАСТЬ ЗАКРЫТА ЗДЕСЬ (проверено 2026-08-08): arch pairing = -ccmWREntry, внедиагональ D0PstarSourceArchOffDiagonalCCMWRCrosswalk.lean:267, диагональ D0PstarSourceArchDiagonalCCMWRCrosswalk.lean:389. tau = ccmW02Entry - ccmWREntry - ccmPrimeEntryN1 (CCMFiniteWeilSourceMatrixN1.lean:97-100), то есть WR входит со знаком минус, а arch pairing равен -WR — знаки требуют сверки, Codex её ведёт. Порядок вывода компонент: WR (E4C) -> W02 (B3.0G) -> Prime (впереди).\par
\end{lemma}

\begin{lemma}[GOAL057\_\allowbreak{}CONTINUUM\_\allowbreak{}NUMERATOR\_\allowbreak{}BRIDGE / step 10]\label{assembly:goal057-continuum-numerator-bridge-10-e72b40147cd8}
\notready
\textbf{Publication state:} \texttt{READY\_\allowbreak{}WITHOUT\_\allowbreak{}EXACT\_\allowbreak{}DECLARATION\_\allowbreak{}RECEIPT}.\par
B3.0F: подъём формы CCM-WR до матрицы\par
\textbf{Assembly supplier field:} \texttt{закрыто 08.08 18:03 'Close finite CCM-WR form lift'}.\par
\textbf{Open receipt:} READY has no exact supplier name and source path.\par
\end{lemma}

\begin{lemma}[GOAL057\_\allowbreak{}CONTINUUM\_\allowbreak{}NUMERATOR\_\allowbreak{}BRIDGE / step 11]\label{assembly:goal057-continuum-numerator-bridge-11-9ca2ec46af65}
\lean{Q3.RouteB.D0Pstar.sourceW02ModePairing_eq_ccmW02Entry}
\leanok
\textbf{Publication state:} \texttt{GREEN}.\par
B3.0G: источниковое W02-спаривание\par
\textbf{Assembly supplier field:} \texttt{sourceW02ModePairing\_\allowbreak{}eq\_\allowbreak{}ccmW02Entry}.\par
\mbox{}\par\smallskip
\begin{Verbatim}[breaklines=true,breakanywhere=true,fontsize=\small]
theorem sourceW02ModePairing_eq_ccmW02Entry (i : PairIndex) (n r : ℤ) :
    sourceW02ModePairing i n r =
      (Q3.RouteB.ccmW02Entry (L_m i) n r : ℂ)
\end{Verbatim}
\smallskip
\end{lemma}
\begin{proof}
\leanok
The proof term is checked in the declaration named above. This receipt does not strengthen its statement.
\end{proof}

\begin{lemma}[GOAL057\_\allowbreak{}CONTINUUM\_\allowbreak{}NUMERATOR\_\allowbreak{}BRIDGE / step 12]\label{assembly:goal057-continuum-numerator-bridge-12-7ac058455310}
\notready
\textbf{Publication state:} \texttt{OPEN\_\allowbreak{}MATH}.\par
обитатель CCMLemma73PreAnchorPort selectedFerrersPreAnchorData (НЕ forall D — C04/\allowbreak{}C10-kill): P.convergence ТЕОРЕМОЙ (бумага 2511.22755)\par
\textbf{Open receipt:} assembly status GAP.\par
\textbf{Recorded note:} РАЗЛОЖЕН НА 9 ЭТАЖЕЙ (вердикт FLOORS 20.08, REQ-D): L73.0 provenance (ждёт Codex, 1/\allowbreak{}10) → L73.1 sourceScale из нормировки I4h0-I0h4, НЕ из вывода (3/\allowbreak{}10) → L73.2 Lemma-7.2 rate O(λ\textasciicircum{}-2) — ГЛАВНАЯ СТЕНА (9/\allowbreak{}10, Meixner-Schaefke Satz 9) → L73.3 оконная ошибка E* (3/\allowbreak{}10) → L73.4 хвост нуль-продолжения ЯВНЫМ ЧЛЕНОМ (4/\allowbreak{}10) → L73.5 Mellin(E*h)=centeredXi — НЕЗАВИСИМ ОТ CODEX, бить параллельно (7/\allowbreak{}10, частично готов) → L73.6 внешний хвост равномерно по подполосе (5/\allowbreak{}10) → L73.7 сборка оценки (4/\allowbreak{}10) → L73.8 запись порта (2/\allowbreak{}10, ПОСЛЕДНЕЙ). MINIMAL\_\allowbreak{}MISSING\_\allowbreak{}IDENTITY: SELECTED\_\allowbreak{}FERRERS\_\allowbreak{}SCALED\_\allowbreak{}GWIN\_\allowbreak{}ERROR\_\allowbreak{}DECOMPOSITION. Счётчик 14/\allowbreak{}14 без изменений. УТОЧНЕНИЕ L73.5 (kill REQ-E, dfe6be5c): цель БЕЗ скаляра УБИТА — Mellin(E*h) = (1/\allowbreak{}4)·centeredXi ТОЧНО (K1: z=0, ratio 0.25 на 79 знаках); новая дыра EXPLICIT\_\allowbreak{}CCM\_\allowbreak{}LIMIT\_\allowbreak{}MELLIN\_\allowbreak{}TO\_\allowbreak{}QUARTER\_\allowbreak{}CENTERED\_\allowbreak{}XI\_\allowbreak{}LEAN; скаляр 4 уходит в sourceScale (L73.1), centeredXi НЕ трогать. СТЕНА L73.2 РАЗЛОЖЕНА (REQ-F, 835d7e97): 7 этажей F72.0-F72.6, ДВА аналитических ядра — F72.1 Satz-9 fixed-mode rate (8/\allowbreak{}10) и F72.3 Fourier-eigenvalue defect |1-chi|<=C/\allowbreak{}λ² (7/\allowbreak{}10, Fuchs 1964); остальное сборка. F72.0 словарь READY\_\allowbreak{}AFTER\_\allowbreak{}OWNER\_\allowbreak{}RELAY\_\allowbreak{}PUSH (ждёт пуша пакета D). Скаляр: portSourceScale = 4·lemma72Scale, четвёрка ровно один раз. F72.0 РЕМОНТ (REQ-G, bc65b407): пакет D РАТИФИЦИРОВАН судьёй (5 аудитов PASS). F72.0A = параметрический замок (n=2j, γ=2π(k+2), γ²=mode4JacobiG — Codex-директива готова, файл G6N1SelectedFerrersPaperParameterDictionary.lean); F72.0B = литеральная привязка SELECTED\_\allowbreak{}FERRERS\_\allowbreak{}TO\_\allowbreak{}SATZ9\_\allowbreak{}REPRESENTATIVE (R1 bind 5/\allowbreak{}10 vs R2 existential-scalar 7/\allowbreak{}10); ЛОВУШКА: def paperPS := наша мода = C10-kill. Скаляр: a\_\allowbreak{}\{j,k\} = h\_\allowbreak{}\{2j,λ\}(0)/\allowbreak{}f\_\allowbreak{}\{j,k\}(0), четвёрка ТОЛЬКО в F72.5. F72.3 SCOPE-LOCK (REQ-H, 3abb8613): crosswalk ДОКАЗАН — a=sqrt(2pi)*lambda, Fuchs c = наша gamma; sqrt(2pi)*chi0=mu\_\allowbreak{}0, sqrt(2pi)*chi2=mu\_\allowbreak{}4, Lambda\_\allowbreak{}n=chi\textasciicircum{}2. Проектно: 1-chi0 \textasciitilde{} 2sqrt2*pi*lambda*e\textasciicircum{}(-4pi lambda\textasciicircum{}2); 1-chi2 \textasciitilde{} (2\textasciicircum{}14/\allowbreak{}3)sqrt2*pi\textasciicircum{}5*lambda\textasciicircum{}9*e\textasciicircum{}(-4pi lambda\textasciicircum{}2) — ЭКСПОНЕНТА, слабый выход C*lambda\textasciicircum{}-2 даром. F72.3 упал 7/\allowbreak{}10 -> 1+3+2 и стал ПАРАЛЛЕЛЕН F72.1. Главная стена теперь одна: F72.1 Meixner-Schaefke Satz 9 (книги на полке НЕТ). ЭТАЖИ ЗАКРЫТЫ ЯДРОМ 20.08 (Linux-тело за Codex): F72.0A параметрический словарь — G6N1SelectedFerrersPaperParameterDictionary.lean, blob 10e9e972, q3\_\allowbreak{}check ok, бонусом gamma = mode4SlepianC (не плодим второй параметр). F72.3A сплетение операторов — G6N1FuchsProjectOperatorIntertwining.lean, blob 59610d91: F\_\allowbreak{}a(U h) = sqrt(2pi)*U(T\_\allowbreak{}lambda h) при единственной гипотезе 0<=lambda. ОСТАЛОСЬ в F72.3: карта собственных чисел Lambda\_\allowbreak{}n = chi\_\allowbreak{}n\textasciicircum{}2 (нужен концентрационный оператор). ГЛАВНАЯ СТЕНА без изменений: F72.1 Meixner-Schaefke Satz 9, книги на полке НЕТ. Строка 12 остаётся КРАСНОЙ: обитатель порта не построен. ИСТОЧНИК СТЕНЫ ДОБЫТ 20.08: Meixner-Schaefke Satz 9, печатная 243 (владелец, Uni-доступ). Равномерно на [-1,1]: ps\_\allowbreak{}n\textasciicircum{}m(z;gamma\textasciicircum{}2) = C(gamma)*(...)*D\_\allowbreak{}\{n-m\}((2gamma)\textasciicircum{}\{1/\allowbreak{}2\}z) + O(gamma\textasciicircum{}-1); D — параболический цилиндр, НЕ наш Эрмит (мост классический, но не проверен). Рейт СХОДИТСЯ с CCM Lemma 7.2: gamma\_\allowbreak{}MS = 2pi*lambda\textasciicircum{}2 => O(gamma\textasciicircum{}-1) = O(lambda\textasciicircum{}-2). ЦЕНА 8/\allowbreak{}10 НЕ УПАЛА — снят только блокер «формулировки нет». НЕ ПРОЧИТАНО: префактор C(gamma), OCR не берёт — это нормировка, читать с чистой страницы. Карточка MEIXNER\_\allowbreak{}SCHAEFKE\_\allowbreak{}1954\_\allowbreak{}USAGE\_\allowbreak{}CARDS.md ПОПРАВКА + ВЕРДИКТ REQ-I (082421cf): сырой остаток Satz 9 = O(gamma\textasciicircum{}-3/\allowbreak{}4), НЕ gamma\textasciicircum{}-1 (моя карточка ошиблась МЕСТОМ числа; судья: SUPERSEDED\_\allowbreak{}RATE\_\allowbreak{}PLACEMENT). Цепочка: raw O(g\textasciicircum{}-3/\allowbreak{}4) -> после деления на (4g/\allowbreak{}pi)\textasciicircum{}(1/\allowbreak{}4) -> O(g\textasciicircum{}-1) -> физически O(lambda\textasciicircum{}-2) = рейт CCM. Префактор ПРОЧИТАН с картинки: (4gamma/\allowbreak{}pi)\textasciicircum{}(1/\allowbreak{}4)/\allowbreak{}sqrt((2n+1)*n!), физические скаляры (1/\allowbreak{}sqrt2)lambda\textasciicircum{}-1/\allowbreak{}2 и (3/\allowbreak{}sqrt2)lambda\textasciicircum{}-1/\allowbreak{}2. ЛОВУШКА: порядок цилиндра nu=n-m=(q-1)/\allowbreak{}2, то есть D\_\allowbreak{}0 и D\_\allowbreak{}4, НЕ D\_\allowbreak{}q (судья убил заранее). ЦЕНА УПАЛА: F72.1 8/\allowbreak{}10 -> 4/\allowbreak{}10 после F72.0B; связка 6/\allowbreak{}10. РАСЩЕПЛЕНИЕ: F72.0B предшественник; F72.1A бумага закрыта/\allowbreak{}порт открыт; F72.1B D->Эрмит LEAN\_\allowbreak{}READY 2/\allowbreak{}10; F72.1C сборка. F72.1B ЗАКРЫТ ЯДРОМ 20.08: G6N1ParabolicCylinderD0D4Exact.lean, blob 5253ef9a, 3 раунда, q3\_\allowbreak{}check ok. D\_\allowbreak{}0 и D\_\allowbreak{}4 при sqrt(4pi)x посчитаны точно через Mathlib Polynomial.hermite (вероятностный, БЕЗ лишней 2\textasciicircum{}(-n/\allowbreak{}2)). СВЕРХ ДИРЕКТИВЫ: explicitCCMLimitH = (1/\allowbreak{}16)D\_\allowbreak{}4 - (3/\allowbreak{}16)D\_\allowbreak{}0 — коэффициент фиксируется членом x\textasciicircum{}4, а член x\textasciicircum{}2 и свободный сходятся САМИ. Нормировка h\_\allowbreak{}0/\allowbreak{}h\_\allowbreak{}4 НЕ вводилась: это F72.0B, объект бумаги нельзя определять нашей формулой. РАЗВИЛКА РЕШЕНА (REQ-J, 7893786e): выбран R2 center-anchored direct cylinder rate, цена ПЕРЕОЦЕНЕНА 7/\allowbreak{}10 -> 5/\allowbreak{}10 ИЗ-ЗА нашего разложения пакета; R1 стала дороже (6/\allowbreak{}10, OVERBUILDS\_\allowbreak{}PUBLIC\_\allowbreak{}OBJECT\_\allowbreak{}GRAPH). h\_\allowbreak{}0/\allowbreak{}h\_\allowbreak{}4 как публичные объекты НЕ НУЖНЫ (FIXED\_\allowbreak{}HERMITE\_\allowbreak{}H0\_\allowbreak{}H4\_\allowbreak{}PUBLIC\_\allowbreak{}OBJECTS\_\allowbreak{}NOT\_\allowbreak{}LOAD\_\allowbreak{}BEARING). Скаляры: a0=1/\allowbreak{}f0(k,0), a4=3/\allowbreak{}f4(k,0) — фиксируются центрами целей D\_\allowbreak{}0(0)=1, D\_\allowbreak{}4(0)=3. F72.0B1 ЗАКРЫТ ЯДРОМ (G6N1CenterAnchorScalarLock.lean, blob a4195440, НОЛЬ ошибок с первого прогона, q3\_\allowbreak{}check ok). Осталось: F72.0B2 source bind 3/\allowbreak{}10 (несущий шов), F72.1A порт 3/\allowbreak{}10, F72.1C сборка 2/\allowbreak{}10, F72.4 2/\allowbreak{}10, F72.5 2/\allowbreak{}10. ВЕРДИКТ REQ-K (a95420e6): выбран S2 (приватный источник-свидетель) + S1 как ОБЩИЙ ПРИЁМНИК ядра; публичный ps НЕ экспортируется, экспортируется центр-нормированное равенство. W13.11 (производная в нуле) ЗАКРЫТ ЯДРОМ: G6N1EvenCenterDerivative.lean, blob 28682c3a, ноль ошибок, q3\_\allowbreak{}check ok — публичная обёртка hasDerivAt\_\allowbreak{}zero\_\allowbreak{}of\_\allowbreak{}even плюс применение к обеим модам плюс приёмник center\_\allowbreak{}data\_\allowbreak{}of\_\allowbreak{}even, который про наши моды НЕ упоминает. ОСТАЛОСЬ несущее в F72.0B2: W13.1 сырой источник независимо от проекта, W13.7 равенство разделяющих собственных значений (ОТДЕЛЬНОЕ обязательство, из G НЕ следует), W13.3/\allowbreak{}W13.4 Lean-порт бумажных фактов, W13.8 подъём координат. ПРИЁМНИК ЗАКРЫТ ЯДРОМ 20.08: G6N1CenterNormalizedUniquenessReceiver.lean, blob 24fdca87, ноль ошибок, q3\_\allowbreak{}check ok. Закрывает СРАЗУ ТРИ LEAN\_\allowbreak{}READY-пункта: F72.0B2C центр-нормированные обёртки, generic\_\allowbreak{}receiver, endpoint\_\allowbreak{}extension. Ключевое свойство — чего в формулировке НЕТ: ни selected Ferrers, ни ps, ни бумажного объекта; f и g оба АРГУМЕНТЫ. Поэтому суррогатное отождествление здесь невозможно по построению. Содержание: деление на скаляр сохраняет однородное линейное уравнение + чётность убивает производную в центре. Продление на концы через closure\_\allowbreak{}Ioo при непрерывности. F72.0B2D КОМПОЗИЦИЯ ЗАКРЫТА ЯДРОМ: G6N1Satz9SourcePackageInterface.lean, blob a7407963, ноль ошибок. Satz9SourceData и ProjectModeData — типизированные дыры; bind доказан БЕЗУСЛОВНО в пакетах: satz9\_\allowbreak{}source\_\allowbreak{}bind\_\allowbreak{}pointwise даёт f(x)/\allowbreak{}f(0) = p(x)/\allowbreak{}p(0) на ЗАМКНУТОМ окне. Ни одно поле Satz9SourceData не упоминает проектных объектов — обитаемо только предъявлением функции и доказательством фактов о ней, переименованием нашей моды НЕ обитаемо. theta несётся ПОЛЕМ, а не выводится: судья пометил theta\_\allowbreak{}equality\_\allowbreak{}is\_\allowbreak{}automatic\_\allowbreak{}from\_\allowbreak{}G=false. ЭТАЖ F72.0B2B ОСТАЁТСЯ ОТКРЫТЫМ — обитателя нет и постулировать его запрещено (абстрактная структура, обитаемая допущением полей, не есть доказательство). ВЕРДИКТ REQ-L (e9c1e470): СУЩЕСТВОВАНИЯ ДОСТАТОЧНО — конструктивная реализация ps НЕ требуется, Classical.choice может выбрать приватного свидетеля, нового глобального аксиома НЕ вводить. Цена остатка: 1/\allowbreak{}10 бумажный аудит + W13.7. НО СУДЬЯ ОПРОВЕРГ МОЁ УТВЕРЖДЕНИЕ: Satz9SourceData есть receiver payload, а НЕ firewall происхождения — отсутствие проектных полей не мешает записать p := наша мода, происхождение есть свойство ЗАВИСИМОСТЕЙ терма, а не имён полей. Математика не затронута (source\_\allowbreak{}interface\_\allowbreak{}theorems\_\allowbreak{}invalidated=false), докстринг исправлен. Починенный дискриминатор: сперва source-only теорема даёт моду И ЕЁ СОБСТВЕННОЕ бумажное собственное значение при gamma\_\allowbreak{}k\textasciicircum{}2, ТОЛЬКО ПОТОМ отдельная теорема доказывает theta\_\allowbreak{}src=theta\_\allowbreak{}project. Экзистенциал, заявленный сразу при theta\_\allowbreak{}project, — kill. СЛЕДУЮЩИЙ НЕСУЩИЙ: W13.7 crosswalk собственных значений. СУЩЕСТВОВАНИЕ НАЙДЕНО В КНИГЕ: Meixner-Schaefke §3.22 Satz 1, печатная 235 — для каждого вещественного gamma\textasciicircum{}2 счётно много собственных значений, к каждому решение непрерывное при z=+-1 вида (1-z\textasciicircum{}2)\textasciicircum{}\{m/\allowbreak{}2\}g(z) с ЦЕЛОЙ g; все собственные значения ВЕЩЕСТВЕННЫ И ПРОСТЫ; ветви lambda\_\allowbreak{}n\textasciicircum{}m(gamma\textasciicircum{}2) с меткой lambda\_\allowbreak{}n\textasciicircum{}m(0)=n(n+1); чётность решения по чётности n-m, при m=0 и n=0,4 ОБЕ ЧЁТНЫЕ. Простота — несущее слово: она даёт единственность с точностью до скаляра, ровно вход приёмника. НЕНУЛЕВОЙ ЦЕНТР СНЯТ КАК ПРОБЕЛ: не цитируется, а ВЫВОДИТСЯ (G6N1EvenSolutionCenterNonvanishing.lean, blob f353622e, ноль ошибок): чётная + значение 0 в центре => совпадает с нулевой функцией по единственности => тождественный ноль; нетривиальная не может. ОСТАЛОСЬ: W13.7 кроссвок ветвей и W13.8/\allowbreak{}9 физический подъём. ВЕРДИКТ REQ-N (d7e6f060): мой третий путь ЗАКОНЕН, но маршрут ГИБРИДНЫЙ — уравнение от DLMF, регулярный спектр от книги; DLMF\_\allowbreak{}30\_\allowbreak{}4\_\allowbreak{}COMPLETENESS\_\allowbreak{}ROUTE НЕ ТРЕБУЕТСЯ. ПРОГНОЗ P\_\allowbreak{}M\_\allowbreak{}NEXT\_\allowbreak{}1 ОПРОВЕРГНУТ (судья, PREDICTION\_\allowbreak{}FATE, retroactive\_\allowbreak{}repair=false): он предсказывал при 0.86, что точный интерфейс DLMF 30.3.5 подтвердит ИСЧЕРПЫВАЮЩИЙ чётный набор решений. Не подтвердил — DLMF пишет односторонне has the solutions. Сбылся ДРУГОЙ, отдельно зарегистрированный исход: наиболее вероятный первый отказ M\_\allowbreak{}REGISTERED\_\allowbreak{}MOST\_\allowbreak{}LIKELY\_\allowbreak{}ONE\_\allowbreak{}WAY\_\allowbreak{}FAILURE = CONFIRMED. Исчерпывающесть заявляет только книга (определённый артикль плюс sämtlich reell und einfach). Моя запись 21.08 называла прогноз подтверждённым — это была ретроактивная починка журнала прогнозов, запрещённая явно; исправлено тем же днём. W13.7D АВТОРИЗОВАН И ЗАКРЫТ ЯДРОМ: G6N1OrderedEnumerationLock.lean, blob fcbdd188, 2 раунда (упал только импорт ℝ). Утверждение АБСТРАКТНОЕ: две строго возрастающие последовательности, чьи части ниже общей отсечки перечисляют одно множество, совпадают почленно, пока обе ниже. Ни сфероидальных функций, ни собственных значений, ни параметра. Продакшн R=2: ординалы 0,1,2 против степеней 0,2,4; средний не потребляется пакетом, но несёт аргумент порядка. ОСТАЛОСЬ: W13.7B порт книжной исчерпывающести в проектный характеристический диапазон, W13.7E перенос пакета. АУДИТ CODEX 21.08 (первый независимый за трое суток): 10/\allowbreak{}10 механически зелёные, 5/\allowbreak{}10 семантических расхождений. ГЛАВНОЕ: приёмник и оба пакета требовали ГЛОБАЛЬНУЮ Continuous вместо ContinuousOn на Icc — а наши моды суть Icc.indicator нуль-продолжения с НЕнулевым значением на конце, значит гипотеза не сильнее, а ЛОЖНА, и endpoint extension был ПУСТ. Ядро не могло увидеть: импликация верна, посылка недостижима. ИСПРАВЛЕНО через Set.EqOn.of\_\allowbreak{}subset\_\allowbreak{}closure, поле переименовано в normalized\_\allowbreak{}continuousOn. Плюс: нетривиальность НЕ следует из полей Satz9SourceData (нулевая функция им удовлетворяет); замок порядка НЕ отождествляет ветви сам. Новые квитанции: receiver 7ec1470f, interface 7c3b781e, nonvanishing db7b2a7d, ordered ce6c9d2e. Гейт-артефакт LINUX\_\allowbreak{}GATE\_\allowbreak{}FIVE\_\allowbreak{}NODES\_\allowbreak{}AFTER\_\allowbreak{}CODEX\_\allowbreak{}AUDIT\_\allowbreak{}2026-08-21.md. ВЕРДИКТ REQ-O (72db5052, ревизия гранта Codex): архитектура трёх тел РАТИФИЦИРОВАНА, грант как написан ОТКЛОНЁН, активация ЗАПРЕЩЕНА, директива HOLD\_\allowbreak{}CODEX\_\allowbreak{}GRANT\_\allowbreak{}REQ\_\allowbreak{}2026\_\allowbreak{}08\_\allowbreak{}21\_\allowbreak{}O. Инвариант: право на пуш и право считать утверждение закрытым — РАЗНЫЕ права. Три статуса SOURCE\_\allowbreak{}WRITTEN -> KERNEL\_\allowbreak{}GREEN -> SEMANTICALLY\_\allowbreak{}ADMITTED; потреблять можно только третий; карантин MAX\_\allowbreak{}KERNEL\_\allowbreak{}GREEN\_\allowbreak{}AWAITING\_\allowbreak{}SEMANTIC\_\allowbreak{}REVIEW=1. Убито: kernel-зелёный как семантический гейт, resume --last как идентичность сессии, pgrep как замок, ID без request blob, три состояния без IN\_\allowbreak{}REVIEW, бессрочный грант, два policy kernel. Ремонт = ОДНА транзакция CODEX\_\allowbreak{}CONTROL\_\allowbreak{}V9\_\allowbreak{}THREE\_\allowbreak{}BODY\_\allowbreak{}LEASE\_\allowbreak{}AND\_\allowbreak{}SEMANTIC\_\allowbreak{}QUARANTINE (версия 8->9, AGENTS.md не трогать) плюс восемь растений; дешевейший решающий тест UNINHABITED\_\allowbreak{}ANTECEDENT\_\allowbreak{}REPLAY. Прогнозы P\_\allowbreak{}O\_\allowbreak{}1..P\_\allowbreak{}O\_\allowbreak{}4 зарегистрированы, не проверены. ВЕРДИКТ REQ-P (2794daa0): control v9 ЗАМОРОЖЕН как at-most-once, владелец может выдать bounded lease. G5: эквивалентность принята (LEAN\_\allowbreak{}PROVED\_\allowbreak{}RESTATEMENT), SelectedCentralFloor и SelectedAnchorRatioData УБРАНЫ из входов G5 (только отсюда!). МОЙ АРГУМЕНТ УБИТ ПО C01: из 1 <= Leak НЕ следует, что масса на краю — нет информации о расположении; фальсификатор — ядра Фейера с быстро растущей степенью. Общая теорема о нормированных рядах ЛОЖНА (постоянная мода). Честная цель одна: source-specific локализация пролатного ряда. Маршрут: источник -> огибающая -> Галёркин при фикс. m -> ОДИН диагональный PairCofinal. Запасной: SelectedLocallyBoundedOnCenteredCriticalStrip (10/\allowbreak{}10 против 9/\allowbreak{}10, цена 7/\allowbreak{}10 против 6/\allowbreak{}10). Отказ ждать у краевой огибающей при sigma->1/\allowbreak{}2. ВЕРДИКТ REQ-Q (17720ac1) — ОБЪЕДИНЕНИЕ ФРОНТОВ: G5 INPUT\_\allowbreak{}A есть EXISTING\_\allowbreak{}GAP\_\allowbreak{}REUSE, не новая стена; переиспользует цепочку F72.1 /\allowbreak{} F72.3 /\allowbreak{} L73.2 /\allowbreak{} L73.3 /\allowbreak{} L73.4 /\allowbreak{} L73.6. Судья ВЫВЕЛ оценку при delta\_\allowbreak{}k <= C*lambda\textasciicircum{}-2 (темп Леммы 7.2): взвешенный интеграл <= C*(lambda\textasciicircum{}(-1/\allowbreak{}2+sigma)/\allowbreak{}(1/\allowbreak{}2+sigma) + lambda\textasciicircum{}-1/\allowbreak{}(1/\allowbreak{}2-sigma)), к нулю при каждом фикс. sigma<1/\allowbreak{}2, НЕ равномерно до 1/\allowbreak{}2. R1 цена 4/\allowbreak{}10 ПОСЛЕ L73.2, убойная сила 10/\allowbreak{}10. СЛЕДСТВИЕ: обитатель BookRegularEvenSpectrum несущий для ДВУХ опор (G6 и G5). Разбухание C\_\allowbreak{}sigma при sigma->1/\allowbreak{}2 РАЗРЕШЕНО; отказ при ОДНОМ фикс. sigma<1/\allowbreak{}2 убивает моментное представление G5, но не Route B; суженный диапазон НЕ достаточен. Переключаться на запасной маршрут НЕТ — он делит ту же стену или требует L73.7. ВЕРДИКТ REQ-S (cab2b8c7): урожай Aristotle РАТИФИЦИРОВАН как partial accept, source-pure, карантин. Мой вопрос про separation+local\_\allowbreak{}finite => infinity УБИТ контрпримером \{0,6\} — конечно, разделено, локально конечно, не бесконечно. Дыра ПЕРЕИМЕНОВАНА: не абстрактная бесконечность, а равномерная high-mode оценка exists N C, C>=0, forall n>=N exists Lambda, RegEvenEig G Lambda and |Lambda-specD(G,n)|<=C, C СНАРУЖИ квантора по n. Выбрана дорога R1 high-mode Jacobi fixed point поверх существующего Spectrum.lean (10/\allowbreak{}10, цена 6/\allowbreak{}10) против R2 min-max/\allowbreak{}Courant требующего ДОКАЗАТЬ compact-resolvent (9/\allowbreak{}10, не следствие уже полученного) и R3 порт книги §1.5 (QUARANTINED, 10/\allowbreak{}10 цена). Codex разрешён READ\_\allowbreak{}ONLY\_\allowbreak{}PLUS\_\allowbreak{}SOURCE\_\allowbreak{}PURE\_\allowbreak{}LOCAL\_\allowbreak{}PROOF\_\allowbreak{}SEARCH в карантинной копии, без Q3-импорта, без правки BookRegularEvenSpectrum. Второй Aristotle ЗАПРЕЩЁН (C09) — грант был на один прогон, потрачен. Частичный обитатель НЕ заводим (C04+C10+W9). Порядок адаптера: закрыть high-mode witness -> even-only пакет -> DLMF-кроссвок -> интерфейс.\par
\end{lemma}

\begin{lemma}[GOAL057\_\allowbreak{}CONTINUUM\_\allowbreak{}NUMERATOR\_\allowbreak{}BRIDGE / step 13]\label{assembly:goal057-continuum-numerator-bridge-13-4306b88f1f13}
\notready
\textbf{Publication state:} \texttt{OPEN\_\allowbreak{}MATH}.\par
равномерная нижняя оценка нормы пробной функции (даёт SelectedTrialNormalizerBounded)\par
\textbf{Open receipt:} assembly status OWNER\_\allowbreak{}DATA.\par
\textbf{Recorded note:} ТОЧНАЯ ФОРМУЛИРОВКА, установлено чтением 2026-08-08. sTrial\_\allowbreak{}m\_\allowbreak{}N = ||gTrial\_\allowbreak{}m\_\allowbreak{}N||\textasciicircum{}\{-1\} (D0KTrialStage3.lean:39). TrialNonzero даёт ПОТОЧЕЧНО: forall k, 0 < ||gTrial\_\allowbreak{}m\_\allowbreak{}N(k)||. Для SelectedTrialNormalizerBounded нужно РАВНОМЕРНО: exists c > 0, forall k, c <= ||gTrial\_\allowbreak{}m\_\allowbreak{}N (selectedPairIndex S k)||. Разница — КВАНТОР, та же болезнь, что с квантором по N: поточечная положительность не даёт отделённости от нуля. Если ||gTrial|| -> 0, то нормализатор -> бесконечность, и произведение ||нормализатор|| * ||хвост|| не сходится, хотя хвост -> 0. Равномерной нижней оценки ||gTrial\_\allowbreak{}m\_\allowbreak{}N|| в проекте НЕТ: grep по le\_\allowbreak{}norm/\allowbreak{}norm\_\allowbreak{}le/\allowbreak{}bddBelow/\allowbreak{}iInf/\allowbreak{}lower/\allowbreak{}uniform пуст. | РАЗБОР 2026-08-08: тоже условие на ДАННЫЕ, не теорема. В структуре ProlateKTrialSourceData УЖЕ ЕСТЬ поле trialNonzero : forall i, TrialNonzero ... (D0ProlateKTrialSource.lean:50) — ПОТОЧЕЧНАЯ ненулевость как требование к источнику. Нужна равномерная версия ТОГО ЖЕ ЖАНРА, соседним полем: trialNormBddBelow : exists c > 0, forall i, c <= ||gTrial\_\allowbreak{}m\_\allowbreak{}N i (prolateCombination (pair i)) (eStar\_\allowbreak{}memLp i)||. Доказывать нечего — надо усилить контракт. | ПРИЁМ B3.0N ПРИМЕРЕН 2026-08-09 — НЕ ПОДХОДИТ. Там оценивают sourceArchimedeanMultiplier t: у него ЯВНАЯ аналитическая формула через digamma, оценка идёт через известную оценку остатка Стилтьеса, ||z|| >= 1/\allowbreak{}4 и логарифмы — всё вычислимо руками. Здесь объект другой: gTrial\_\allowbreak{}m\_\allowbreak{}N i h hLp = P\_\allowbreak{}m\_\allowbreak{}N i (gTrial\_\allowbreak{}m i h hLp) (D0KTrialStage2.lean:56-62) — НОРМА ПРОЕКЦИИ элемента, приходящего из данных (prolateCombination (pair i)). Явной формулы нет. Норма проекции может быть сколь угодно мала, если элемент почти ортогонален подпространству E\_\allowbreak{}m\_\allowbreak{}N — никакая аналитика этого не запретит. Подтверждает: шаг 13 действительно OWNER\_\allowbreak{}DATA, а не 'приём ещё не найден'. | 20.08: для N2 НЕ нужен (zero-mode сокращение нормировщиков); остаётся входом N2-этажей как контракт\par
\end{lemma}

\begin{lemma}[GOAL057\_\allowbreak{}CONTINUUM\_\allowbreak{}NUMERATOR\_\allowbreak{}BRIDGE / step 14]\label{assembly:goal057-continuum-numerator-bridge-14-eab96e1fe42d}
\notready
\textbf{Publication state:} \texttt{OPEN\_\allowbreak{}MATH}.\par
N2: SelectedNormalizedGalerkinMellinCompactDecay (compact-open, НЕ Hilbert-norm)\par
\textbf{Assembly supplier field:} \texttt{selectedProjectionTailDecay\_\allowbreak{}of\_\allowbreak{}physicalFourierEnergyControl}.\par
\textbf{Open receipt:} assembly status GAP.\par
\textbf{Recorded note:} разложена на этажи N2\_\allowbreak{}0..N2\_\allowbreak{}5 вердиктом 20.08; MINIMAL\_\allowbreak{}MISSING\_\allowbreak{}IDENTITY = source-scaled Mellin projection tail rate; исполнение ДО обитателей D и порта запрещено\par
\end{lemma}

\begin{lemma}[GOAL057\_\allowbreak{}CONTINUUM\_\allowbreak{}NUMERATOR\_\allowbreak{}BRIDGE / step 15]\label{assembly:goal057-continuum-numerator-bridge-15-3e283e9f905f}
\notready
\textbf{Publication state:} \texttt{OPEN\_\allowbreak{}MATH}.\par
SelectedPhysicalFourierEnergyControl — суммируемость и ограниченность энергий\par
\textbf{Open receipt:} assembly status OWNER\_\allowbreak{}DATA.\par
\textbf{Recorded note:} Предикат-конъюнкция, объявлен PFE:66. (а) для каждого k ряд sum\_\allowbreak{}n physicalFourierWeight i n * ||physicalFourierCoefficient i (gTrial\_\allowbreak{}m ...) n||\textasciicircum{}2 суммируем; (б) IsBoundedUnder (<=) atTop (norm . selectedPhysicalFourierEnergy S) — энергии ограничены ПО k. Часть (б) — снова равномерность, см. правило pointwise-vs-uniform. Поставщика нет: объявлен и потреблён в одном файле. | РАЗБОР 2026-08-08: тот же жанр. Энергии считаются от gTrial\_\allowbreak{}m, который из S.source.pair — то есть от ДАННЫХ. Часть (а) суммируемость и часть (б) ограниченность по k обе зависят от выбора pair. Место для полей — рядом с eStar\_\allowbreak{}memLp (D0ProlateKTrialSource.lean:46), где уже стоят требования того же вида.\par
\end{lemma}

\begin{lemma}[GOAL057\_\allowbreak{}CONTINUUM\_\allowbreak{}NUMERATOR\_\allowbreak{}BRIDGE / step 16]\label{assembly:goal057-continuum-numerator-bridge-16-d8effc5abb58}
\notready
\textbf{Publication state:} \texttt{OPEN\_\allowbreak{}MATH}.\par
SelectedPhysicalBandwidthCofinal — полоса уходит в бесконечность\par
\textbf{Open receipt:} assembly status OWNER\_\allowbreak{}DATA.\par
\textbf{Recorded note:} Предикат: Tendsto (fun k => physicalFourierBandwidth (selectedPairIndex S k)) atTop atTop, объявлен PFE:81. Утверждение о ВЫБОРЕ пути (selectedPairIndex), а не о математике — кандидат на дешёвое закрытие: может следовать из конструкции выбора. Поставщика нет. Проверять первым из трёх. | РАЗБОР 2026-08-08: НЕДОКАЗУЕМО, это условие на ДАННЫЕ. physicalFourierBandwidth i = 2*pi*(N+1) /\allowbreak{} log(m) (PFE:47 + logLength). selectedPairIndex берётся из полей структуры CanonicalData: parent, extract — произвольные функции, не формула. В CanonicalData УЖЕ ЕСТЬ поле parentCofinal : PairCofinal parent, требующее m -> oo И N -> oo (D0CanonicalApproximation.lean:66-71). НО ЭТОГО НЕ ХВАТАЕТ: при N\_\allowbreak{}k = k и m\_\allowbreak{}k = exp(k\textasciicircum{}2) оба стремятся к бесконечности, а 2*pi*(k+1)/\allowbreak{}k\textasciicircum{}2 -> 0. Нужно ДОПОЛНИТЕЛЬНОЕ соотношение скоростей: N\_\allowbreak{}k /\allowbreak{} log(m\_\allowbreak{}k) -> oo. Место для поля — рядом с parentCofinal в CanonicalData.\par
\end{lemma}

\begin{lemma}[GOAL057\_\allowbreak{}CONTINUUM\_\allowbreak{}NUMERATOR\_\allowbreak{}BRIDGE / step 17]\label{assembly:goal057-continuum-numerator-bridge-17-10385b07305d}
\notready
\textbf{Publication state:} \texttt{READY\_\allowbreak{}WITHOUT\_\allowbreak{}EXACT\_\allowbreak{}DECLARATION\_\allowbreak{}RECEIPT}.\par
B3.0H: подъём W02 до конечной формы\par
\textbf{Assembly supplier field:} \texttt{D0PstarSourceW02FiniteFormCCMW02Crosswalk}.\par
\textbf{Open receipt:} no unique exact name in aristotle\_\allowbreak{}proofs.db.\par
\textbf{Recorded note:} закрыто 09.08 00:02\par
\end{lemma}

\begin{lemma}[GOAL057\_\allowbreak{}CONTINUUM\_\allowbreak{}NUMERATOR\_\allowbreak{}BRIDGE / step 18]\label{assembly:goal057-continuum-numerator-bridge-18-5e0f22efa496}
\notready
\textbf{Publication state:} \texttt{READY\_\allowbreak{}WITHOUT\_\allowbreak{}EXACT\_\allowbreak{}DECLARATION\_\allowbreak{}RECEIPT}.\par
B3.0I: источниковое prime-спаривание\par
\textbf{Assembly supplier field:} \texttt{D0PstarSourcePrimeModePairing}.\par
\textbf{Open receipt:} no unique exact name in aristotle\_\allowbreak{}proofs.db.\par
\textbf{Recorded note:} закрыто 09.08 00:58\par
\end{lemma}

\begin{lemma}[GOAL057\_\allowbreak{}CONTINUUM\_\allowbreak{}NUMERATOR\_\allowbreak{}BRIDGE / step 19]\label{assembly:goal057-continuum-numerator-bridge-19-753d92932964}
\notready
\textbf{Publication state:} \texttt{READY\_\allowbreak{}WITHOUT\_\allowbreak{}EXACT\_\allowbreak{}DECLARATION\_\allowbreak{}RECEIPT}.\par
B3.0J: подъём prime до конечной формы\par
\textbf{Assembly supplier field:} \texttt{D0PstarSourcePrimeFiniteFormCCMPrimeCrosswalk}.\par
\textbf{Open receipt:} no unique exact name in aristotle\_\allowbreak{}proofs.db.\par
\textbf{Recorded note:} закрыто 09.08 01:27\par
\end{lemma}

\begin{lemma}[GOAL057\_\allowbreak{}CONTINUUM\_\allowbreak{}NUMERATOR\_\allowbreak{}BRIDGE / step 20]\label{assembly:goal057-continuum-numerator-bridge-20-d137d2112525}
\lean{Q3.RouteB.D0Pstar.sourceWeilFiniteForm_eq_ccmWeilMatrixForm}
\leanok
\textbf{Publication state:} \texttt{GREEN}.\par
B3.0K: ПОЛНАЯ форма Вейля = матричная форма CCM\par
\textbf{Assembly supplier field:} \texttt{sourceWeilFiniteForm\_\allowbreak{}eq\_\allowbreak{}ccmWeilMatrixForm}.\par
\mbox{}\par\smallskip
\begin{Verbatim}[breaklines=true,breakanywhere=true,fontsize=\small]
theorem sourceWeilFiniteForm_eq_ccmWeilMatrixForm (i : PairIndex)
    (c d : CCMModeFinite i.N → ℂ) :
    ((∑ j, ∑ k,
        star (c j) *
          sourceW02ModePairing i
            (ccmModeFinite i.N j)
            (ccmModeFinite i.N k) *
          d k
\end{Verbatim}
\smallskip
\end{lemma}
\begin{proof}
\leanok
The proof term is checked in the declaration named above. This receipt does not strengthen its statement.
\end{proof}

\begin{lemma}[GOAL057\_\allowbreak{}CONTINUUM\_\allowbreak{}NUMERATOR\_\allowbreak{}BRIDGE / step 21]\label{assembly:goal057-continuum-numerator-bridge-21-0492d6e1ed0e}
\notready
\textbf{Publication state:} \texttt{READY\_\allowbreak{}WITHOUT\_\allowbreak{}EXACT\_\allowbreak{}DECLARATION\_\allowbreak{}RECEIPT}.\par
B3.0L: изометрия Фурье в L2\par
\textbf{Assembly supplier field:} \texttt{D0PstarSourceLogWindowFourierL2Isometry}.\par
\textbf{Open receipt:} no unique exact name in aristotle\_\allowbreak{}proofs.db.\par
\textbf{Recorded note:} закрыто 09.08 03:54\par
\end{lemma}

\begin{lemma}[GOAL057\_\allowbreak{}CONTINUUM\_\allowbreak{}NUMERATOR\_\allowbreak{}BRIDGE / step 22]\label{assembly:goal057-continuum-numerator-bridge-22-05855ed748ed}
\notready
\textbf{Publication state:} \texttt{READY\_\allowbreak{}WITHOUT\_\allowbreak{}EXACT\_\allowbreak{}DECLARATION\_\allowbreak{}RECEIPT}.\par
B3.0M: конечный Фурье-реестр\par
\textbf{Assembly supplier field:} \texttt{D0PstarSourceWeilFiniteFourierLedger}.\par
\textbf{Open receipt:} no unique exact name in aristotle\_\allowbreak{}proofs.db.\par
\textbf{Recorded note:} закрыто 09.08 05:08\par
\end{lemma}

\begin{lemma}[GOAL057\_\allowbreak{}CONTINUUM\_\allowbreak{}NUMERATOR\_\allowbreak{}BRIDGE / step 23]\label{assembly:goal057-continuum-numerator-bridge-23-8ff79dfa8a53}
\lean{Q3.RouteB.D0Pstar.sourceArchimedeanMultiplier_add_explicitShift_nonneg}
\leanok
\textbf{Publication state:} \texttt{GREEN}.\par
B3.0N: РАВНОМЕРНАЯ нижняя оценка арх. символа\par
\textbf{Assembly supplier field:} \texttt{sourceArchimedeanMultiplier\_\allowbreak{}add\_\allowbreak{}explicitShift\_\allowbreak{}nonneg}.\par
\mbox{}\par\smallskip
\begin{Verbatim}[breaklines=true,breakanywhere=true,fontsize=\small]
theorem sourceArchimedeanMultiplier_add_explicitShift_nonneg (t : ℝ) :
    0 ≤ sourceArchimedeanMultiplier t +
      (|Real.log Real.pi| + Real.log 4 + 6)
\end{Verbatim}
\smallskip
\end{lemma}
\begin{proof}
\leanok
The proof term is checked in the declaration named above. This receipt does not strengthen its statement.
\end{proof}

\begin{lemma}[GOAL057\_\allowbreak{}CONTINUUM\_\allowbreak{}NUMERATOR\_\allowbreak{}BRIDGE / step 24]\label{assembly:goal057-continuum-numerator-bridge-24-fe66a1c2ab8a}
\notready
\textbf{Publication state:} \texttt{READY\_\allowbreak{}WITHOUT\_\allowbreak{}EXACT\_\allowbreak{}DECLARATION\_\allowbreak{}RECEIPT}.\par
B3.0O: сдвинутый арх. sqrt-вес\par
\textbf{Assembly supplier field:} \texttt{D0PstarShiftedArchSqrtWeight}.\par
\textbf{Open receipt:} no unique exact name in aristotle\_\allowbreak{}proofs.db.\par
\textbf{Recorded note:} закрыто 09.08 07:20\par
\end{lemma}

\begin{lemma}[GOAL057\_\allowbreak{}CONTINUUM\_\allowbreak{}NUMERATOR\_\allowbreak{}BRIDGE / step 25]\label{assembly:goal057-continuum-numerator-bridge-25-c5b2d9f5b30e}
\notready
\textbf{Publication state:} \texttt{READY\_\allowbreak{}WITHOUT\_\allowbreak{}EXACT\_\allowbreak{}DECLARATION\_\allowbreak{}RECEIPT}.\par
обитатель SelectedProlatePreAnchorData (пакет: index,pair,кофинальности,lambda\_\allowbreak{}eq,MemLp)\par
\textbf{Assembly supplier field:} \texttt{selectedFerrersPreAnchorData (def, kernel-green 20.08, Linux-тело за Codex)}.\par
\textbf{Open receipt:} no unique exact name in aristotle\_\allowbreak{}proofs.db.\par
\textbf{Recorded note:} Расписание precommit k↦(k+2,k+2,K=5(k+2)); пара = точный свидетель exists\_\allowbreak{}modeZero\_\allowbreak{}modeFour\_\allowbreak{}selectedFerrersProductionProlatePair; pair\_\allowbreak{}spec экспортирован (требование судьи выполнено). Профиль стандартная тройка, q3\_\allowbreak{}check ok. SOURCE RECORD: docs/\allowbreak{}routeB\_\allowbreak{}bus/\allowbreak{}LINUX\_\allowbreak{}SOURCE\_\allowbreak{}RECORD\_\allowbreak{}SELECTED\_\allowbreak{}FERRERS\_\allowbreak{}PREANCHOR\_\allowbreak{}DATA\_\allowbreak{}2026-08-20.md\par
\end{lemma}

\begin{lemma}[GOAL057\_\allowbreak{}CONTINUUM\_\allowbreak{}NUMERATOR\_\allowbreak{}BRIDGE / step 26]\label{assembly:goal057-continuum-numerator-bridge-26-5ea822d1d70b}
\notready
\textbf{Publication state:} \texttt{OPEN\_\allowbreak{}MATH}.\par
N3: same-family locally-uniform crosswalk D0Pstar->Muntz\par
\textbf{Open receipt:} assembly status GAP.\par
\textbf{Recorded note:} сборка (вердикт G6-узлов 19.08); комбинирует N1-композер с N2\par
\end{lemma}

\begin{lemma}[GOAL057\_\allowbreak{}CONTINUUM\_\allowbreak{}NUMERATOR\_\allowbreak{}BRIDGE / step 27]\label{assembly:goal057-continuum-numerator-bridge-27-89bb1222d93e}
\notready
\textbf{Publication state:} \texttt{OPEN\_\allowbreak{}MATH}.\par
N4: SlotS2 из фиксированного selected-предела\par
\textbf{Open receipt:} assembly status GAP.\par
\textbf{Recorded note:} строгая потребительская сборка; квантификация по всем ClusterData\par
\end{lemma}

\section{Pillar G2 — finite-cell validation package}

\begin{lemma}[PSD\_\allowbreak{}CERTIFICATE\_\allowbreak{}FOR\_\allowbreak{}CCM\_\allowbreak{}CELL / step 1]\label{assembly:psd-certificate-for-ccm-cell-1-8cf53f076eff}
\lean{Q3.RouteB.ccmWeilMatFinite_transpose_eq}
\leanok
\textbf{Publication state:} \texttt{GREEN}.\par
K эрмитова (симметричная)\par
\textbf{Assembly supplier field:} \texttt{ccmWeilMatFinite\_\allowbreak{}transpose\_\allowbreak{}eq}.\par
\mbox{}\par\smallskip
\begin{Verbatim}[breaklines=true,breakanywhere=true,fontsize=\small]
theorem ccmWeilMatFinite_transpose_eq (mProject N : ℕ)
    (hm : 2 ≤ mProject)
    (hN : 1 ≤ N) :
    (ccmWeilMatFinite mProject N).transpose =
      ccmWeilMatFinite mProject N
\end{Verbatim}
\smallskip
\end{lemma}
\begin{proof}
\leanok
The proof term is checked in the declaration named above. This receipt does not strengthen its statement.
\end{proof}

\begin{lemma}[PSD\_\allowbreak{}CERTIFICATE\_\allowbreak{}FOR\_\allowbreak{}CCM\_\allowbreak{}CELL / step 2]\label{assembly:psd-certificate-for-ccm-cell-2-6f5e84df52b3}
\notready
\textbf{Publication state:} \texttt{READY\_\allowbreak{}WITHOUT\_\allowbreak{}EXACT\_\allowbreak{}DECLARATION\_\allowbreak{}RECEIPT}.\par
поблочная симметрия tau = W02 - WR - Prime\par
\textbf{Assembly supplier field:} \texttt{ccmW02Entry\_\allowbreak{}symm + ccmWREntry\_\allowbreak{}symm + ccmPrimeEntryN1\_\allowbreak{}symm}.\par
\textbf{Open receipt:} no unique exact name in aristotle\_\allowbreak{}proofs.db.\par
\textbf{Recorded note:} каждый блок доказан отдельно: 291, 343, 310; итог ccmWeilTauN1\_\allowbreak{}symm:381\par
\end{lemma}

\begin{lemma}[PSD\_\allowbreak{}CERTIFICATE\_\allowbreak{}FOR\_\allowbreak{}CCM\_\allowbreak{}CELL / step 3]\label{assembly:psd-certificate-for-ccm-cell-3-850c4346aef8}
\lean{Q3.RouteB.ccmWeilMatFinite_centrosymmetric}
\leanok
\textbf{Publication state:} \texttt{GREEN}.\par
J инволюция с JKJ = K\par
\textbf{Assembly supplier field:} \texttt{ccmWeilMatFinite\_\allowbreak{}centrosymmetric}.\par
\mbox{}\par\smallskip
\begin{Verbatim}[breaklines=true,breakanywhere=true,fontsize=\small]
theorem ccmWeilMatFinite_centrosymmetric (mProject N : ℕ)
    (hm : 2 ≤ mProject)
    (hN : 1 ≤ N)
    (i j : CCMModeFinite N) :
    ccmWeilMatFinite mProject N
        (ccmNegFinite N i) (ccmNegFinite N j) =
      ccmWeilMatFinite mProject N i j
\end{Verbatim}
\smallskip
\end{lemma}
\begin{proof}
\leanok
The proof term is checked in the declaration named above. This receipt does not strengthen its statement.
\end{proof}

\begin{lemma}[PSD\_\allowbreak{}CERTIFICATE\_\allowbreak{}FOR\_\allowbreak{}CCM\_\allowbreak{}CELL / step 4]\label{assembly:psd-certificate-for-ccm-cell-4-382189942164}
\lean{Q3.RouteB.ccmWeilOpFinite_commutes_reflection}
\leanok
\textbf{Publication state:} \texttt{GREEN}.\par
оператор коммутирует с отражением\par
\textbf{Assembly supplier field:} \texttt{ccmWeilOpFinite\_\allowbreak{}commutes\_\allowbreak{}reflection}.\par
\mbox{}\par\smallskip
\begin{Verbatim}[breaklines=true,breakanywhere=true,fontsize=\small]
theorem ccmWeilOpFinite_commutes_reflection (mProject N : ℕ) (hm : 2 ≤ mProject) (hN : 1 ≤ N) :
    (ccmReflectionEndFinite N).comp (ccmWeilOpFinite mProject N) =
      (ccmWeilOpFinite mProject N).comp (ccmReflectionEndFinite N)
\end{Verbatim}
\smallskip
\end{lemma}
\begin{proof}
\leanok
The proof term is checked in the declaration named above. This receipt does not strengthen its statement.
\end{proof}

\begin{lemma}[PSD\_\allowbreak{}CERTIFICATE\_\allowbreak{}FOR\_\allowbreak{}CCM\_\allowbreak{}CELL / step 5]\label{assembly:psd-certificate-for-ccm-cell-5-5dd91daead72}
\lean{Q3.RouteB.realPosDef_map_complex}
\leanok
\textbf{Publication state:} \texttt{GREEN}.\par
переход R -> C\par
\textbf{Assembly supplier field:} \texttt{realPosDef\_\allowbreak{}map\_\allowbreak{}complex}.\par
\mbox{}\par\smallskip
\begin{Verbatim}[breaklines=true,breakanywhere=true,fontsize=\small]
theorem realPosDef_map_complex {n : Type*} [Fintype n] [DecidableEq n]
    (Q : Matrix n n ℝ) (hQ : Q.PosDef) :
    (Q.map (algebraMap ℝ ℂ)).PosDef
\end{Verbatim}
\smallskip
\end{lemma}
\begin{proof}
\leanok
The proof term is checked in the declaration named above. This receipt does not strengthen its statement.
\end{proof}

\begin{lemma}[PSD\_\allowbreak{}CERTIFICATE\_\allowbreak{}FOR\_\allowbreak{}CCM\_\allowbreak{}CELL / step 6]\label{assembly:psd-certificate-for-ccm-cell-6-4455ed5ab925}
\notready
\textbf{Publication state:} \texttt{READY\_\allowbreak{}WITHOUT\_\allowbreak{}EXACT\_\allowbreak{}DECLARATION\_\allowbreak{}RECEIPT}.\par
G PosDef\par
\textbf{Assembly supplier field:} \texttt{Matrix.PosDef.one (Mathlib)}.\par
\textbf{Open receipt:} no unique exact name in aristotle\_\allowbreak{}proofs.db.\par
\textbf{Recorded note:} подтверждено тремя способами: (1) hbottom использует x .. x без весовой матрицы; (2) CCMModeFinite = Fin (2N+1), метрики нет; (3) ccmWeilOpFinite = mulVecLin, без метрики. Остаток: прямой леммы Matrix.PosDef 1 в Mathlib не нашёл, доказывается в две строки | Статус READY -> GAP: note сам признаёт «прямой леммы Matrix.PosDef 1 в Mathlib не нашёл, доказывается в две строки». «Доказывается» != «доказано». Две строки — но их никто не написал. | ЗАКРЫТО 2026-08-08: лемма СУЩЕСТВУЕТ в Mathlib — Matrix.PosDef.one, PosDef.lean:204, доказана двумя строками (nontriviality R; .diagonal fun i => zero\_\allowbreak{}lt\_\allowbreak{}one' R). Прежний note «прямой леммы не нашёл» неверен: искали по имени posDef\_\allowbreak{}one, а она называется PosDef.one. Требования: StarOrderedRing R, DecidableEq n, NoZeroDivisors R — все выполнены для C.\par
\end{lemma}

\begin{lemma}[PSD\_\allowbreak{}CERTIFICATE\_\allowbreak{}FOR\_\allowbreak{}CCM\_\allowbreak{}CELL / step 7]\label{assembly:psd-certificate-for-ccm-cell-7-3fbf05d27a48}
\lean{Q3.RouteB.ccmShiftedWeilMatFinite_posSemidef_of_bottomRayleigh}
\leanok
\textbf{Publication state:} \texttt{GREEN}.\par
сертификат K - beta*G + tau*(Gq)(Gq)\textasciicircum{}H >= 0\par
\textbf{Assembly supplier field:} \texttt{ccmShiftedWeilMatFinite\_\allowbreak{}posSemidef\_\allowbreak{}of\_\allowbreak{}bottomRayleigh}.\par
\mbox{}\par\smallskip
\begin{Verbatim}[breaklines=true,breakanywhere=true,fontsize=\small]
theorem ccmShiftedWeilMatFinite_posSemidef_of_bottomRayleigh (mProject N : ℕ) (epsilon : ℝ)
    (hm : 2 ≤ mProject) (hN : 1 ≤ N)
    (hbottom : ∀ x : CCMModeFinite N → ℝ,
      epsilon * (x ⬝ᵥ x) ≤
        x ⬝ᵥ Matrix.mulVec (ccmWeilMatFinite mProject N) x) :
    (ccmShifted
\end{Verbatim}
\smallskip
\end{lemma}
\begin{proof}
\leanok
The proof term is checked in the declaration named above. This receipt does not strengthen its statement.
\end{proof}

\begin{validation}[PSD\_\allowbreak{}CERTIFICATE\_\allowbreak{}FOR\_\allowbreak{}CCM\_\allowbreak{}CELL / step 8]\label{assembly:psd-certificate-for-ccm-cell-8-90f874f53042}
\notready
\textbf{Publication state:} \texttt{VALIDATION\_\allowbreak{}ONLY}.\par
рэлеевская нижняя граница epsilon\par
\textbf{Assembly supplier field:} \texttt{Arb interval LDL @ 240 dps, все 121+120 пивотов положительны}.\par
\textbf{Open receipt:} not a Lean proof authority.\par
\textbf{Recorded note:} CCM\_\allowbreak{}CONTROL\_\allowbreak{}CELL\_\allowbreak{}CERT\_\allowbreak{}INTERVAL\_\allowbreak{}PASS; на 120 dps INSUFFICIENT\_\allowbreak{}PRECISION (интервалы шире величины), на 240 dps INTERVAL\_\allowbreak{}POSITIVE\_\allowbreak{}DEFINITE; повтор 180->360 подтвердил | ЧИСЛО, НЕ ТЕОРЕМА. Источник — JSON с интервальными вычислениями Arb, не объект Lean. Статус изменён READY -> NUMERIC: до подъёма измерения до формального PosSemidef цепочка формально не закрыта. Механизм подъёма существует: PSD\_\allowbreak{}PenaltyCertificate.lean (рациональный сертификат, 0 sorry) + Matrix.posSemidef\_\allowbreak{}iff\_\allowbreak{}dotProduct\_\allowbreak{}mulVec. К этой матрице не применён. | 2026-08-20: снят с критического пути крыши вердиктом REQ-2026-08-20-B (клетка = falsifier, не поставщик SlotH2a)\par
\end{validation}

\begin{lemma}[PSD\_\allowbreak{}CERTIFICATE\_\allowbreak{}FOR\_\allowbreak{}CCM\_\allowbreak{}CELL / step 9]\label{assembly:psd-certificate-for-ccm-cell-9-929fc080301d}
\lean{Q3.RouteB.ccmShiftedWeil_rankOneCorrection_kernel_and_weightedSymmetric}
\leanok
\textbf{Publication state:} \texttt{GREEN}.\par
ранг-один поправка tau*(Gq)(Gq)\textasciicircum{}H\par
\textbf{Assembly supplier field:} \texttt{ccmShiftedWeil\_\allowbreak{}rankOneCorrection\_\allowbreak{}kernel\_\allowbreak{}and\_\allowbreak{}weightedSymmetric}.\par
\mbox{}\par\smallskip
\begin{Verbatim}[breaklines=true,breakanywhere=true,fontsize=\small]
theorem ccmShiftedWeil_rankOneCorrection_kernel_and_weightedSymmetric (mProject N : ℕ) (epsilon : ℝ)
    (xi : CCMModeFinite N → ℝ)
    (hm : 2 ≤ mProject) (hN : 1 ≤ N)
    (heig : Matrix.mulVec (ccmWeilMatFinite mProject N) xi = epsilon • xi)
    (hxiEven : ∀ i, xi (ccmNegFinite N i
\end{Verbatim}
\smallskip
\end{lemma}
\begin{proof}
\leanok
The proof term is checked in the declaration named above. This receipt does not strengthen its statement.
\end{proof}

\begin{lemma}[PSD\_\allowbreak{}CERTIFICATE\_\allowbreak{}FOR\_\allowbreak{}CCM\_\allowbreak{}CELL / step 10]\label{assembly:psd-certificate-for-ccm-cell-10-c64458c81598}
\lean{Q3.RouteB.exists_ccmEta_normalized_even_eigenvector_of_simple_even_eigenvector}
\leanok
\textbf{Publication state:} \texttt{GREEN}.\par
q: J-чётный, нормированный eta*q = 1\par
\textbf{Assembly supplier field:} \texttt{exists\_\allowbreak{}ccmEta\_\allowbreak{}normalized\_\allowbreak{}even\_\allowbreak{}eigenvector\_\allowbreak{}of\_\allowbreak{}simple\_\allowbreak{}even\_\allowbreak{}eigenvector}.\par
\mbox{}\par\smallskip
\begin{Verbatim}[breaklines=true,breakanywhere=true,fontsize=\small]
theorem exists_ccmEta_normalized_even_eigenvector_of_simple_even_eigenvector (mProject N : ℕ)
    (epsilon : ℝ)
    (xi : CCMModeFinite N → ℝ)
    (hm : 2 ≤ mProject)
    (hN : 1 ≤ N)
    (hxi0 : xi ≠ 0)
    (heig :
      Matrix.mulVec (ccmWeilMatFinite mProject N) xi = epsilon • xi)
    (h
\end{Verbatim}
\smallskip
\end{lemma}
\begin{proof}
\leanok
The proof term is checked in the declaration named above. This receipt does not strengthen its statement.
\end{proof}

\begin{lemma}[PSD\_\allowbreak{}CERTIFICATE\_\allowbreak{}FOR\_\allowbreak{}CCM\_\allowbreak{}CELL / step 11]\label{assembly:psd-certificate-for-ccm-cell-11-837ee0658815}
\lean{Q3.RouteB.ccmEigenvector_even_of_simple_eigenspace_and_normalized}
\leanok
\textbf{Publication state:} \texttt{GREEN}.\par
чётность собственного вектора\par
\textbf{Assembly supplier field:} \texttt{ccmEigenvector\_\allowbreak{}even\_\allowbreak{}of\_\allowbreak{}simple\_\allowbreak{}eigenspace\_\allowbreak{}and\_\allowbreak{}normalized}.\par
\mbox{}\par\smallskip
\begin{Verbatim}[breaklines=true,breakanywhere=true,fontsize=\small]
theorem ccmEigenvector_even_of_simple_eigenspace_and_normalized (mProject N : ℕ) (epsilon : ℝ)
    (xi : CCMModeFinite N → ℝ)
    (hm : 2 ≤ mProject) (hN : 1 ≤ N)
    (heig : Matrix.mulVec (ccmWeilMatFinite mProject N) xi =
      epsilon • xi)
    (hnormalized : ccmEtaFinite N 
\end{Verbatim}
\smallskip
\end{lemma}
\begin{proof}
\leanok
The proof term is checked in the declaration named above. This receipt does not strengthen its statement.
\end{proof}

\begin{lemma}[PSD\_\allowbreak{}CERTIFICATE\_\allowbreak{}FOR\_\allowbreak{}CCM\_\allowbreak{}CELL / step 12]\label{assembly:psd-certificate-for-ccm-cell-12-4919d31e7c4b}
\lean{Q3.RouteB.ccmEtaFinite_dotProduct_ne_zero_of_even_simple_eigenvector}
\leanok
\textbf{Publication state:} \texttt{GREEN}.\par
eta*xi != 0 (нормируемость)\par
\textbf{Assembly supplier field:} \texttt{ccmEtaFinite\_\allowbreak{}dotProduct\_\allowbreak{}ne\_\allowbreak{}zero\_\allowbreak{}of\_\allowbreak{}even\_\allowbreak{}simple\_\allowbreak{}eigenvector}.\par
\mbox{}\par\smallskip
\begin{Verbatim}[breaklines=true,breakanywhere=true,fontsize=\small]
theorem ccmEtaFinite_dotProduct_ne_zero_of_even_simple_eigenvector (mProject N : ℕ)
    (epsilon : ℝ)
    (xi : CCMModeFinite N → ℝ)
    (hm : 2 ≤ mProject)
    (hN : 1 ≤ N)
    (hxi0 : xi ≠ 0)
    (heig :
      Matrix.mulVec (ccmWeilMatFinite mProject N) xi = epsilon • xi)
    (h
\end{Verbatim}
\smallskip
\end{lemma}
\begin{proof}
\leanok
The proof term is checked in the declaration named above. This receipt does not strengthen its statement.
\end{proof}

\begin{validation}[PSD\_\allowbreak{}CERTIFICATE\_\allowbreak{}FOR\_\allowbreak{}CCM\_\allowbreak{}CELL / step 13]\label{assembly:psd-certificate-for-ccm-cell-13-47e858572416}
\notready
\textbf{Publication state:} \texttt{VALIDATION\_\allowbreak{}ONLY}.\par
a = <q,Kq>\par
\textbf{Assembly supplier field:} \texttt{a = 4.71998e-59 (строгий интервал Arb)}.\par
\textbf{Open receipt:} not a Lean proof authority.\par
\textbf{Recorded note:} вычислено для J-чётной проекции q, как требует контракт; прежние 5.37e-59 были для комплексного source-вектора | ЧИСЛО, НЕ ТЕОРЕМА. Источник — JSON с интервальными вычислениями Arb, не объект Lean. Статус изменён READY -> NUMERIC: до подъёма измерения до формального PosSemidef цепочка формально не закрыта. Механизм подъёма существует: PSD\_\allowbreak{}PenaltyCertificate.lean (рациональный сертификат, 0 sorry) + Matrix.posSemidef\_\allowbreak{}iff\_\allowbreak{}dotProduct\_\allowbreak{}mulVec. К этой матрице не применён. | 2026-08-20: снят с критического пути крыши вердиктом REQ-2026-08-20-B (клетка = falsifier, не поставщик SlotH2a)\par
\end{validation}

\begin{validation}[PSD\_\allowbreak{}CERTIFICATE\_\allowbreak{}FOR\_\allowbreak{}CCM\_\allowbreak{}CELL / step 14]\label{assembly:psd-certificate-for-ccm-cell-14-95ad6a9d255b}
\notready
\textbf{Publication state:} \texttt{VALIDATION\_\allowbreak{}ONLY}.\par
beta с a < beta\par
\textbf{Assembly supplier field:} \texttt{beta = 1e-56, beta - a = 9.95280e-57 > 0}.\par
\textbf{Open receipt:} not a Lean proof authority.\par
\textbf{Recorded note:} tau = 1, без оптимизации; Phase 2 ищет beta*\_\allowbreak{}N = min(odd floor, q-perp compression floor) | ЧИСЛО, НЕ ТЕОРЕМА. Источник — JSON с интервальными вычислениями Arb, не объект Lean. Статус изменён READY -> NUMERIC: до подъёма измерения до формального PosSemidef цепочка формально не закрыта. Механизм подъёма существует: PSD\_\allowbreak{}PenaltyCertificate.lean (рациональный сертификат, 0 sorry) + Matrix.posSemidef\_\allowbreak{}iff\_\allowbreak{}dotProduct\_\allowbreak{}mulVec. К этой матрице не применён. | 2026-08-20: снят с критического пути крыши вердиктом REQ-2026-08-20-B (клетка = falsifier, не поставщик SlotH2a)\par
\end{validation}

\begin{lemma}[PSD\_\allowbreak{}CERTIFICATE\_\allowbreak{}FOR\_\allowbreak{}CCM\_\allowbreak{}CELL / step 15]\label{assembly:psd-certificate-for-ccm-cell-15-a2fb5939faa5}
\lean{H2aPenalty.simplicity_clause}
\leanok
\textbf{Publication state:} \texttt{GREEN}.\par
простота нижнего собственного подпространства\par
\textbf{Assembly supplier field:} \texttt{simplicity\_\allowbreak{}clause}.\par
\mbox{}\par\smallskip
\begin{Verbatim}[breaklines=true,breakanywhere=true,fontsize=\small]
theorem simplicity_clause {G K : Matrix n n ℂ} {q : n → ℂ} {β τ lam : ℝ}
    (hG : G.PosDef)
    (hcert : (K - (β : ℂ) • G + (τ : ℂ) • Matrix.vecMulVec (G *ᵥ q) (star (G *ᵥ q))).PosSemidef)
    (hlamβ : lam < β) :
    ∀ x y, GEig K G
\end{Verbatim}
\smallskip
\end{lemma}
\begin{proof}
\leanok
The proof term is checked in the declaration named above. This receipt does not strengthen its statement.
\end{proof}

\section{Goal 058 — ground diagonal to centered Xi}

\begin{lemma}[REALZERO\_\allowbreak{}GROUND\_\allowbreak{}DIAGONAL\_\allowbreak{}TO\_\allowbreak{}XI / step 0]\label{assembly:realzero-ground-diagonal-to-xi-0-07d9540d1f3b}
\notready
\textbf{Publication state:} \texttt{OPEN\_\allowbreak{}MATH}.\par
точный объект, координата и нормировка зафиксированы\par
\textbf{Assembly supplier field:} \texttt{координаты источника заданы, расписание и нормировка — частично}.\par
\textbf{Open receipt:} assembly status GAP.\par
\textbf{Recorded note:} G0 OPEN\_\allowbreak{}PARTIAL. Целые узлы источника и полюсы 2*pi*k/\allowbreak{}L определены; кофинальное расписание (m\_\allowbreak{}j, N\_\allowbreak{}j) и точная невырожденная нормировка не закреплены.\par
\end{lemma}

\begin{lemma}[REALZERO\_\allowbreak{}GROUND\_\allowbreak{}DIAGONAL\_\allowbreak{}TO\_\allowbreak{}XI / step 1]\label{assembly:realzero-ground-diagonal-to-xi-1-82bcd32d3c3f}
\notready
\textbf{Publication state:} \texttt{OPEN\_\allowbreak{}MATH}.\par
кофинальный конечный simple-even ground-пакет\par
\textbf{Assembly supplier field:} \texttt{penalty-сертификат (A) · блочное расщепление чётности (B) · GLOWER/\allowbreak{}Yoshida (C) · Schur/\allowbreak{}Feshbach (D) · ранг-инерция (E) · внешняя теорема (F)}.\par
\textbf{Open receipt:} assembly status GAP.\par
\textbf{Recorded note:} G1 OPEN\_\allowbreak{}MAIN\_\allowbreak{}SPECTRAL\_\allowbreak{}FRONT. Нужен пакет (epsilon\_\allowbreak{}j, xi\_\allowbreak{}j, heig, hbottom, hsimple, hnormalized) на кофинальном пути. Чётность отдельно НЕ поставляется — её выводит обёртка ...\_\allowbreak{}simple\_\allowbreak{}normalized. Шесть поставщиков-кандидатов, ни один не доведён.\par
\end{lemma}

\begin{lemma}[REALZERO\_\allowbreak{}GROUND\_\allowbreak{}DIAGONAL\_\allowbreak{}TO\_\allowbreak{}XI / step 2]\label{assembly:realzero-ground-diagonal-to-xi-2-8b026aa968c4}
\lean{Q3.RouteB.ccmSourceLagrangePolynomial_complex_zerosRealOn_of_bottomRayleigh_simple_normalized}
\leanok
\textbf{Publication state:} \texttt{GREEN}.\par
вещественность нулей лагранжева многочлена ground-строки\par
\textbf{Assembly supplier field:} \texttt{ccmSourceLagrangePolynomial\_\allowbreak{}complex\_\allowbreak{}zerosRealOn\_\allowbreak{}of\_\allowbreak{}bottomRayleigh\_\allowbreak{}simple\_\allowbreak{}normalized}.\par
\mbox{}\par\smallskip
\begin{Verbatim}[breaklines=true,breakanywhere=true,fontsize=\small]
theorem ccmSourceLagrangePolynomial_complex_zerosRealOn_of_bottomRayleigh_simple_normalized {ι : Type*} [Fintype ι] [DecidableEq ι]
    (mProject N : ℕ) (epsilon : ℝ)
    (xi : CCMModeFinite N → ℝ)
    (hm : 2 ≤ mProject) (hN : 1 ≤ N)
    (heig : Matrix.mulVec (ccmWeilMatFinite mProject N) xi =
 
\end{Verbatim}
\smallskip
\end{lemma}
\begin{proof}
\leanok
The proof term is checked in the declaration named above. This receipt does not strengthen its statement.
\end{proof}

\begin{lemma}[REALZERO\_\allowbreak{}GROUND\_\allowbreak{}DIAGONAL\_\allowbreak{}TO\_\allowbreak{}XI / step 3]\label{assembly:realzero-ground-diagonal-to-xi-3-65d7ba3e02fa}
\lean{Q3.RouteB.Proposition59GroundLagrangeZeroSetBridge}
\leanok
\textbf{Publication state:} \texttt{GREEN}.\par
перенос МНОЖЕСТВА нулей с многочлена на преобразование Proposition-5.9\par
\textbf{Assembly supplier field:} \texttt{Proposition59GroundLagrangeZeroSetBridge}.\par
\mbox{}\par\smallskip
\begin{Verbatim}[breaklines=true,breakanywhere=true,fontsize=\small]
theorem Proposition59GroundLagrangeZeroSetBridge {ι : Type*} [Fintype ι] [DecidableEq ι] (mProject N : ℕ) (epsilon : ℝ) (xi : CCMModeFinite N → ℝ) (hm : 2 ≤ mProject) (hN : 1 ≤ N) (heig : Matrix.mulVec (ccmWeilMatFinite mProject N) xi = epsilon • xi) (hnormalized : ccmEtaFinite N ⬝ᵥ xi = 1) (hbottom : ∀ x : CCMModeFinite N → ℝ, epsilon * (x ⬝ᵥ x) ≤ x ⬝ᵥ Matrix.mulVec (ccmWeilMatFinite mProject N) x) (hsimple : Module.finrank ℝ ((ccmWeilOpFinite mProject N).eigenspace epsilon) = 1) (b : Module.Ba
\end{Verbatim}
\smallskip
\end{lemma}
\begin{proof}
\leanok
The proof term is checked in the declaration named above. This receipt does not strengthen its statement.
\end{proof}

\begin{lemma}[REALZERO\_\allowbreak{}GROUND\_\allowbreak{}DIAGONAL\_\allowbreak{}TO\_\allowbreak{}XI / step 4]\label{assembly:realzero-ground-diagonal-to-xi-4-45464de255f3}
\notready
\textbf{Publication state:} \texttt{OPEN\_\allowbreak{}MATH}.\par
та же F\_\allowbreak{}j отслеживает projected prolate trial локально равномерно\par
\textbf{Assembly supplier field:} \texttt{невязка делить на истинный зазор (A) · Feshbach-граф (B) · penalty-overlap (C) · defect-Gram (D) · low/\allowbreak{}high split (E) · norm-resolvent Galerkin (F)}.\par
\textbf{Open receipt:} assembly status GAP.\par
\textbf{Recorded note:} G3 MAIN\_\allowbreak{}WALL. Главная стена маршрута. Переносится СХОДИМОСТЬ на ground-семью, не вещественность на trial: вещественность неустойчива к возмущению. Убитое решение: exact ground equals trial.\par
\end{lemma}

\begin{lemma}[REALZERO\_\allowbreak{}GROUND\_\allowbreak{}DIAGONAL\_\allowbreak{}TO\_\allowbreak{}XI / step 5]\label{assembly:realzero-ground-diagonal-to-xi-5-3b3cb2299564}
\notready
\textbf{Publication state:} \texttt{OPEN\_\allowbreak{}MATH}.\par
projected trial отслеживает continuum CCM trial\par
\textbf{Assembly supplier field:} \texttt{не назначен}.\par
\textbf{Open receipt:} assembly status GAP.\par
\textbf{Recorded note:} G3c. Проекция конечного среза на континуальный пробник CCM.\par
\end{lemma}

\begin{lemma}[REALZERO\_\allowbreak{}GROUND\_\allowbreak{}DIAGONAL\_\allowbreak{}TO\_\allowbreak{}XI / step 6]\label{assembly:realzero-ground-diagonal-to-xi-6-3b26e681fbfb}
\notready
\textbf{Publication state:} \texttt{OPEN\_\allowbreak{}MATH}.\par
CCM Lemma 7.3: континуальное trial-преобразование сходится к Xi\par
\textbf{Assembly supplier field:} \texttt{CCM Lemma 7.3 — доказана в статье, порт в проект отсутствует}.\par
\textbf{Open receipt:} assembly status GAP.\par
\textbf{Recorded note:} G4 PAPER\_\allowbreak{}PROVED\_\allowbreak{}PROJECT\_\allowbreak{}IMPORT\_\allowbreak{}OPEN. Это тот же объект, что doc-alias hermfact1: статус PAPER\_\allowbreak{}PROVED, Lean-порт OPEN. Локально равномерно на замкнутых подполосах.\par
\end{lemma}

\begin{lemma}[REALZERO\_\allowbreak{}GROUND\_\allowbreak{}DIAGONAL\_\allowbreak{}TO\_\allowbreak{}XI / step 7]\label{assembly:realzero-ground-diagonal-to-xi-7-3369afe5ee81}
\notready
\textbf{Publication state:} \texttt{READY\_\allowbreak{}WITHOUT\_\allowbreak{}EXACT\_\allowbreak{}DECLARATION\_\allowbreak{}RECEIPT}.\par
zero-escape: вещественные нули аппроксимантов запрещают невещественные у предела\par
\textbf{Assembly supplier field:} \texttt{rh\_\allowbreak{}of\_\allowbreak{}canonical\_\allowbreak{}strip\_\allowbreak{}slots · ZerosApproachOn · ClassicalXiInterface}.\par
\textbf{Open receipt:} no unique exact name in aristotle\_\allowbreak{}proofs.db.\par
\textbf{Recorded note:} G5. Логическое ядро ДОКАЗАНО в крыше: одна фиксированная выбранная семья с вещественными нулями и нужным пределом даёт Q3.RH. sameCofinalGuard не позволяет подставить независимую диагональ. Аналитический перенос в крыше доступен.\par
\end{lemma}

\section{Pillar G3 — real-zeros bridge}

\begin{lemma}[SIMPLE\_\allowbreak{}EVEN\_\allowbreak{}GROUND\_\allowbreak{}TO\_\allowbreak{}REAL\_\allowbreak{}ZEROS / step 1]\label{assembly:simple-even-ground-to-real-zeros-1-52b339bd3c39}
\notready
\textbf{Publication state:} \texttt{OPEN\_\allowbreak{}MATH}.\par
вещественность нулей для CCM-семьи\par
\textbf{Assembly supplier field:} \texttt{ccmSourceLagrangePolynomial\_\allowbreak{}complex\_\allowbreak{}zerosRealOn\_\allowbreak{}of\_\allowbreak{}bottomRayleigh\_\allowbreak{}simple}.\par
\textbf{Open receipt:} assembly status MISMATCH.\par
\textbf{Recorded note:} доказана для КОНКРЕТНОЙ семьи; слот сформулирован для абстрактной C.Pstar.family\par
\end{lemma}

\begin{lemma}[SIMPLE\_\allowbreak{}EVEN\_\allowbreak{}GROUND\_\allowbreak{}TO\_\allowbreak{}REAL\_\allowbreak{}ZEROS / step 2]\label{assembly:simple-even-ground-to-real-zeros-2-8476cdef606e}
\notready
\textbf{Publication state:} \texttt{READY\_\allowbreak{}WITHOUT\_\allowbreak{}EXACT\_\allowbreak{}DECLARATION\_\allowbreak{}RECEIPT}.\par
determinant identity\par
\textbf{Assembly supplier field:} \texttt{zerosRealOn\_\allowbreak{}of\_\allowbreak{}hermitian\_\allowbreak{}charpoly\_\allowbreak{}mul (charpoly в факторизации)}.\par
\textbf{Open receipt:} no unique exact name in aristotle\_\allowbreak{}proofs.db.\par
\end{lemma}

\begin{lemma}[SIMPLE\_\allowbreak{}EVEN\_\allowbreak{}GROUND\_\allowbreak{}TO\_\allowbreak{}REAL\_\allowbreak{}ZEROS / step 3]\label{assembly:simple-even-ground-to-real-zeros-3-18f9ba94255c}
\lean{Q3.RouteB.posDefSelfAdjoint_exists_hermitian}
\leanok
\textbf{Publication state:} \texttt{GREEN}.\par
modified-Hilbert self-adjoint descent\par
\textbf{Assembly supplier field:} \texttt{posDefSelfAdjoint\_\allowbreak{}exists\_\allowbreak{}hermitian}.\par
\mbox{}\par\smallskip
\begin{Verbatim}[breaklines=true,breakanywhere=true,fontsize=\small]
theorem posDefSelfAdjoint_exists_hermitian {n : Type*} [Fintype n] [DecidableEq n]
    (Q D : Matrix n n ℂ) (hQ : Q.PosDef)
    (hSA : Q * D = Dᴴ * Q) :
    ∃ H : Matrix n n ℂ, H.IsHermitian ∧ H.charpoly = D.charpoly
\end{Verbatim}
\smallskip
\end{lemma}
\begin{proof}
\leanok
The proof term is checked in the declaration named above. This receipt does not strengthen its statement.
\end{proof}

\begin{lemma}[SIMPLE\_\allowbreak{}EVEN\_\allowbreak{}GROUND\_\allowbreak{}TO\_\allowbreak{}REAL\_\allowbreak{}ZEROS / step 4]\label{assembly:simple-even-ground-to-real-zeros-4-bf14c9f1776c}
\notready
\textbf{Publication state:} \texttt{READY\_\allowbreak{}WITHOUT\_\allowbreak{}EXACT\_\allowbreak{}DECLARATION\_\allowbreak{}RECEIPT}.\par
complement/\allowbreak{}lattice factor\par
\textbf{Assembly supplier field:} \texttt{realFactor с гипотезой ZerosRealOn}.\par
\textbf{Open receipt:} no unique exact name in aristotle\_\allowbreak{}proofs.db.\par
\end{lemma}

\begin{lemma}[SIMPLE\_\allowbreak{}EVEN\_\allowbreak{}GROUND\_\allowbreak{}TO\_\allowbreak{}REAL\_\allowbreak{}ZEROS / step 5]\label{assembly:simple-even-ground-to-real-zeros-5-8f55e2c5fdb5}
\notready
\textbf{Publication state:} \texttt{READY\_\allowbreak{}WITHOUT\_\allowbreak{}EXACT\_\allowbreak{}DECLARATION\_\allowbreak{}RECEIPT}.\par
nonvanishing phase\par
\textbf{Assembly supplier field:} \texttt{hunit: forall z, unit z != 0}.\par
\textbf{Open receipt:} no unique exact name in aristotle\_\allowbreak{}proofs.db.\par
\end{lemma}

\begin{lemma}[SIMPLE\_\allowbreak{}EVEN\_\allowbreak{}GROUND\_\allowbreak{}TO\_\allowbreak{}REAL\_\allowbreak{}ZEROS / step 6]\label{assembly:simple-even-ground-to-real-zeros-6-1d037e5dd3d3}
\notready
\textbf{Publication state:} \texttt{OPEN\_\allowbreak{}MATH}.\par
H2aAt как конкретный предикат\par
\textbf{Open receipt:} assembly status GAP.\par
\textbf{Recorded note:} свободный параметр Index -> Prop; должен упаковать ШЕСТЬ гипотез поставщика: heig, hnormalized, hbottom, hsimple, базис b, hm/\allowbreak{}hN. hxiEven ИСКЛЮЧЁН — выводится, см. шаг 11.\par
\end{lemma}

\begin{lemma}[SIMPLE\_\allowbreak{}EVEN\_\allowbreak{}GROUND\_\allowbreak{}TO\_\allowbreak{}REAL\_\allowbreak{}ZEROS / step 7]\label{assembly:simple-even-ground-to-real-zeros-7-26b419cc4c24}
\notready
\textbf{Publication state:} \texttt{READY\_\allowbreak{}WITHOUT\_\allowbreak{}EXACT\_\allowbreak{}DECLARATION\_\allowbreak{}RECEIPT}.\par
базис фактора по ядру сдвинутой формы\par
\textbf{Assembly supplier field:} \texttt{Module.finBasis (Mathlib) — базис строится, не постулируется}.\par
\textbf{Open receipt:} no unique exact name in aristotle\_\allowbreak{}proofs.db.\par
\textbf{Recorded note:} в гипотезах дан как данное, не как существование; при инстанциации нужен конструктивно или через exists | КАНДИДАТ НА ДЕШЁВОЕ ЗАКРЫТИЕ (2026-08-08, компиляцией НЕ проверено): Module.finBasis (Mathlib/\allowbreak{}LinearAlgebra/\allowbreak{}Dimension/\allowbreak{}Free.lean:230) строит базис конечномерного модуля. Наш фактор (CCMModeFinite N -> R) /\allowbreak{} ker (toBilin' T) конечномерен над R, значит базис не обязан быть гипотезой — его можно построить. Нужны инстансы Module.Finite и Module.Free; для фактора конечномерного по подпространству должны выводиться автоматически. Проверять запуском lean, не глазами. | ПРОВЕРЕНО LEAN 2026-08-08: Module.finBasis R \_\allowbreak{} строит базис фактора (CCMModeFinite N -> R) /\allowbreak{} ker (toBilin' (ccmShiftedWeilMatFinite ...)). lake env lean прошёл без ошибок. Инстансы Module.Finite/\allowbreak{}Module.Free выводятся автоматически. ВНИМАНИЕ: код проверен вне репозитория (WRITE\_\allowbreak{}LOCK, роль reader) — требует внесения Codex-ом, чтобы шаг стал закрытым и в проекте.\par
\end{lemma}

\begin{lemma}[SIMPLE\_\allowbreak{}EVEN\_\allowbreak{}GROUND\_\allowbreak{}TO\_\allowbreak{}REAL\_\allowbreak{}ZEROS / step 8]\label{assembly:simple-even-ground-to-real-zeros-8-ef03e67a9d8e}
\notready
\textbf{Publication state:} \texttt{READY\_\allowbreak{}WITHOUT\_\allowbreak{}EXACT\_\allowbreak{}DECLARATION\_\allowbreak{}RECEIPT}.\par
согласование простоты: eigenspace vs обобщённая задача (K,G)\par
\textbf{Assembly supplier field:} \texttt{finrank\_\allowbreak{}eigenspace\_\allowbreak{}eq\_\allowbreak{}one\_\allowbreak{}of\_\allowbreak{}geig\_\allowbreak{}simple (доказана, вне репы)}.\par
\textbf{Open receipt:} no unique exact name in aristotle\_\allowbreak{}proofs.db.\par
\textbf{Recorded note:} H2a\_\allowbreak{}SimpleEvenGround\_\allowbreak{}FromPenaltyCoercivity (H2aPenaltyCoercivity.lean:395) даёт простоту для ОБОБЩЁННОЙ задачи GEig K G (две матрицы, пучок). Поставщику нужен finrank eigenspace ОБЫЧНОЙ задачи (одна матрица W). Переходника пучок->обычная в проекте НЕТ. ВНИМАНИЕ: не путать с шагом 9 — тот закрывает ker<->eigenspace ВНУТРИ обычной задачи и к этому разрыву отношения не имеет. | МЕХАНИЗМ СВЕДЕНИЯ НАЙДЕН 2026-08-08 (компиляцией не проверен). GEig определён как x != 0 AND K *v x = mu . (G *v x) (H2aPenaltyCoercivity.lean:50). В нашей PSD-цепочке G = единичная матрица (шаг 6 той цепочки, Matrix.PosDef.one). При G = 1 имеем G *v x = x по Matrix.one\_\allowbreak{}mulVec, значит GEig K 1 mu x <=> x != 0 AND K *v x = mu . x — это ОБЫЧНАЯ задача, то есть x in Module.End.eigenspace K.mulVecLin mu. Тогда simplicity\_\allowbreak{}clause (пропорциональность любых двух) + exists\_\allowbreak{}lowest (существование ненулевого) дают finrank eigenspace = 1. Порода: ПЕРЕХОДНИК, цена часы, а не неизвестно. В Mathlib пучка (K,G) НЕТ — genEigenspace это жордановы цепочки (f-mu)\textasciicircum{}k, ложное совпадение по слову generalized. | ПРОВЕРЕНО LEAN 2026-08-08, шаг собран ЦЕЛИКОМ: (1) geig\_\allowbreak{}one\_\allowbreak{}iff — при G=1 GEig K 1 mu x <=> x != 0 AND K *v x = mu . x, через Matrix.one\_\allowbreak{}mulVec; (2) finrank\_\allowbreak{}eigenspace\_\allowbreak{}eq\_\allowbreak{}one\_\allowbreak{}of\_\allowbreak{}geig\_\allowbreak{}simple — из простоты пучка следует finrank (End.eigenspace K.mulVecLin lam) = 1, через finrank\_\allowbreak{}eq\_\allowbreak{}one\_\allowbreak{}iff\_\allowbreak{}of\_\allowbreak{}nonzero' (FiniteDimensional/\allowbreak{}Basic.lean:545). lake env lean прошёл без ошибок. Код вне репозитория (WRITE\_\allowbreak{}LOCK) — требует внесения.\par
\end{lemma}

\begin{lemma}[SIMPLE\_\allowbreak{}EVEN\_\allowbreak{}GROUND\_\allowbreak{}TO\_\allowbreak{}REAL\_\allowbreak{}ZEROS / step 9]\label{assembly:simple-even-ground-to-real-zeros-9-f63d8ac71417}
\lean{Q3.RouteB.zerosRealOn_of_zerosApproachOn}
\leanok
\textbf{Publication state:} \texttt{GREEN}.\par
предельный переход сохраняет вещественность\par
\textbf{Assembly supplier field:} \texttt{zerosRealOn\_\allowbreak{}of\_\allowbreak{}zerosApproachOn}.\par
\mbox{}\par\smallskip
\begin{Verbatim}[breaklines=true,breakanywhere=true,fontsize=\small]
theorem zerosRealOn_of_zerosApproachOn (S : Set ℂ) (F : ℕ → ℂ → ℂ) (f : ℂ → ℂ)
    (hF : ∀ n, ZerosRealOn Set.univ (F n))
    (htransfer : ZerosApproachOn S F f) :
    ZerosRealOn S f
\end{Verbatim}
\smallskip
\end{lemma}
\begin{proof}
\leanok
The proof term is checked in the declaration named above. This receipt does not strengthen its statement.
\end{proof}

\begin{lemma}[SIMPLE\_\allowbreak{}EVEN\_\allowbreak{}GROUND\_\allowbreak{}TO\_\allowbreak{}REAL\_\allowbreak{}ZEROS / step 10]\label{assembly:simple-even-ground-to-real-zeros-10-eedc567aa6d3}
\lean{Q3.RouteB.montel_anchor_nonzero_limit}
\leanok
\textbf{Publication state:} \texttt{GREEN}.\par
конечное -> бесконечное с живым якорем\par
\textbf{Assembly supplier field:} \texttt{montel\_\allowbreak{}anchor\_\allowbreak{}nonzero\_\allowbreak{}limit}.\par
\mbox{}\par\smallskip
\begin{Verbatim}[breaklines=true,breakanywhere=true,fontsize=\small]
theorem montel_anchor_nonzero_limit (f : ℕ → ℂ → ℂ) (a c : ℂ)
    (hf : ∀ n, Differentiable ℂ (f n))
    (hbdd : ∀ K : Set ℂ, IsCompact K → ∃ M : ℝ, ∀ n, ∀ z ∈ K, ‖f n z‖ ≤ M)
    (hA : ∀ n, f n a = c) (hc : c ≠ 0) :
    ∃ e : ℕ → ℕ, StrictM
\end{Verbatim}
\smallskip
\end{lemma}
\begin{proof}
\leanok
The proof term is checked in the declaration named above. This receipt does not strengthen its statement.
\end{proof}

\begin{lemma}[SIMPLE\_\allowbreak{}EVEN\_\allowbreak{}GROUND\_\allowbreak{}TO\_\allowbreak{}REAL\_\allowbreak{}ZEROS / step 11]\label{assembly:simple-even-ground-to-real-zeros-11-3459437f3684}
\lean{Q3.RouteB.rh_iff_centeredXi_zeros_real}
\leanok
\textbf{Publication state:} \texttt{GREEN}.\par
выход в классическую RH\par
\textbf{Assembly supplier field:} \texttt{rh\_\allowbreak{}iff\_\allowbreak{}centeredXi\_\allowbreak{}zeros\_\allowbreak{}real}.\par
\mbox{}\par\smallskip
\begin{Verbatim}[breaklines=true,breakanywhere=true,fontsize=\small]
theorem rh_iff_centeredXi_zeros_real  : Q3.RH ↔ CenteredXiZerosReal
\end{Verbatim}
\smallskip
\end{lemma}
\begin{proof}
\leanok
The proof term is checked in the declaration named above. This receipt does not strengthen its statement.
\end{proof}

\begin{lemma}[SIMPLE\_\allowbreak{}EVEN\_\allowbreak{}GROUND\_\allowbreak{}TO\_\allowbreak{}REAL\_\allowbreak{}ZEROS / step 12]\label{assembly:simple-even-ground-to-real-zeros-12-49fc9156e3f3}
\lean{Q3.RouteB.ccmShiftedWeilOpFinite_ker_eq_eigenspace}
\leanok
\textbf{Publication state:} \texttt{GREEN}.\par
ker сдвинутой формы <-> eigenspace (обычная задача)\par
\textbf{Assembly supplier field:} \texttt{ccmShiftedWeilOpFinite\_\allowbreak{}ker\_\allowbreak{}eq\_\allowbreak{}eigenspace}.\par
\mbox{}\par\smallskip
\begin{Verbatim}[breaklines=true,breakanywhere=true,fontsize=\small]
theorem ccmShiftedWeilOpFinite_ker_eq_eigenspace (mProject N : ℕ) (epsilon : ℝ) :
    LinearMap.ker (ccmShiftedWeilOpFinite mProject N epsilon) =
      (ccmWeilOpFinite mProject N).eigenspace epsilon
\end{Verbatim}
\smallskip
\end{lemma}
\begin{proof}
\leanok
The proof term is checked in the declaration named above. This receipt does not strengthen its statement.
\end{proof}

\begin{lemma}[SIMPLE\_\allowbreak{}EVEN\_\allowbreak{}GROUND\_\allowbreak{}TO\_\allowbreak{}REAL\_\allowbreak{}ZEROS / step 13]\label{assembly:simple-even-ground-to-real-zeros-13-acb96483cd2d}
\lean{Q3.RouteB.ccmShiftedWeilMatFinite_posSemidef_of_bottomRayleigh}
\leanok
\textbf{Publication state:} \texttt{GREEN}.\par
Rayleigh-граница снизу -> PosSemidef сдвинутой матрицы\par
\textbf{Assembly supplier field:} \texttt{ccmShiftedWeilMatFinite\_\allowbreak{}posSemidef\_\allowbreak{}of\_\allowbreak{}bottomRayleigh}.\par
\mbox{}\par\smallskip
\begin{Verbatim}[breaklines=true,breakanywhere=true,fontsize=\small]
theorem ccmShiftedWeilMatFinite_posSemidef_of_bottomRayleigh (mProject N : ℕ) (epsilon : ℝ)
    (hm : 2 ≤ mProject) (hN : 1 ≤ N)
    (hbottom : ∀ x : CCMModeFinite N → ℝ,
      epsilon * (x ⬝ᵥ x) ≤
        x ⬝ᵥ Matrix.mulVec (ccmWeilMatFinite mProject N) x) :
    (ccmShifted
\end{Verbatim}
\smallskip
\end{lemma}
\begin{proof}
\leanok
The proof term is checked in the declaration named above. This receipt does not strengthen its statement.
\end{proof}

\begin{lemma}[SIMPLE\_\allowbreak{}EVEN\_\allowbreak{}GROUND\_\allowbreak{}TO\_\allowbreak{}REAL\_\allowbreak{}ZEROS / step 14]\label{assembly:simple-even-ground-to-real-zeros-14-cc9c79ddcbcd}
\lean{Q3.RouteB.ccmEigenvector_even_of_simple_eigenspace_and_normalized}
\leanok
\textbf{Publication state:} \texttt{GREEN}.\par
чётность собственного вектора (hxiEven)\par
\textbf{Assembly supplier field:} \texttt{ccmEigenvector\_\allowbreak{}even\_\allowbreak{}of\_\allowbreak{}simple\_\allowbreak{}eigenspace\_\allowbreak{}and\_\allowbreak{}normalized}.\par
\mbox{}\par\smallskip
\begin{Verbatim}[breaklines=true,breakanywhere=true,fontsize=\small]
theorem ccmEigenvector_even_of_simple_eigenspace_and_normalized (mProject N : ℕ) (epsilon : ℝ)
    (xi : CCMModeFinite N → ℝ)
    (hm : 2 ≤ mProject) (hN : 1 ≤ N)
    (heig : Matrix.mulVec (ccmWeilMatFinite mProject N) xi =
      epsilon • xi)
    (hnormalized : ccmEtaFinite N 
\end{Verbatim}
\smallskip
\end{lemma}
\begin{proof}
\leanok
The proof term is checked in the declaration named above. This receipt does not strengthen its statement.
\end{proof}

\begin{lemma}[SIMPLE\_\allowbreak{}EVEN\_\allowbreak{}GROUND\_\allowbreak{}TO\_\allowbreak{}REAL\_\allowbreak{}ZEROS / step 15]\label{assembly:simple-even-ground-to-real-zeros-15-5373ac5b9f5f}
\notready
\textbf{Publication state:} \texttt{OPEN\_\allowbreak{}MATH}.\par
нормирующий множитель c\_\allowbreak{}N и сходимость F\_\allowbreak{}N -> centeredXi\par
\textbf{Open receipt:} assembly status GAP.\par
\textbf{Recorded note:} sourceLagrangePolynomial = sum\_\allowbreak{}k xi(k) * prod\_\allowbreak{}\{j!=k\}(lam j - X): при 2N+1 модах произведение из 2N сомножителей порядка N, рост \textasciitilde{}N\textasciicircum{}(2N). centeredXi конечна. Без нормировки c\_\allowbreak{}N сходимость невозможна в принципе. c\_\allowbreak{}N в проекте ОТСУТСТВУЕТ (ccmL — параметр модели; '\_\allowbreak{}normalized' в Parity:161 нормирует СОБСТВЕННЫЙ ВЕКТОР, не полином). Связь sourceLagrangePolynomial <-> centeredXi: ноль теорем. ZeroEscapeLogic.lean:48 сам фиксирует: 'independent of the still-open object'.\par
\end{lemma}

\section{Registered semantic edges}
This compact appendix lists only declarations on exact registered supplier-to-consumer paths; unrelated helper declarations are omitted.
\begin{lemma}[Historical edge record E001]
\textbf{Supplier:} \texttt{Q3.RouteB.D0Pstar.selectedFerrersAbelLogZeroExtension\_\allowbreak{}fourier\_\allowbreak{}decay\_\allowbreak{}quantitative}.\par
\textbf{Consumer:} \texttt{Q3.RouteB.D0Pstar.selectedFerrersAbelLimitHm\_\allowbreak{}physicalCoefficient\_\allowbreak{}le}.\par
\textbf{Recorded historical relation:} \texttt{DIRECT}.\par
\textbf{Current relation:} \texttt{DIRECT}.\par
\textbf{Consumer source state:} \texttt{CURRENT}.\par
\end{lemma}
\begin{lemma}[Historical edge record E002]
\textbf{Supplier:} \texttt{Q3.RouteB.D0Pstar.sourceWeilGraphAmbient\_\allowbreak{}evenNonzeroMode}.\par
\textbf{Consumer:} \texttt{Q3.RouteB.D0Pstar.sourceWeilGraphAmbient\_\allowbreak{}evenNonzeroMode\_\allowbreak{}orthonormal}.\par
\textbf{Recorded historical relation:} \texttt{DIRECT}.\par
\textbf{Current relation:} \texttt{DIRECT}.\par
\textbf{Consumer source state:} \texttt{CURRENT}.\par
\end{lemma}
\begin{lemma}[Historical edge record E003]
\textbf{Supplier:} \texttt{Q3.RouteB.D0Pstar.sourceWeilGraphAmbient\_\allowbreak{}evenZeroMode}.\par
\textbf{Consumer:} \texttt{Q3.RouteB.D0Pstar.sourceWeilGraphAmbient\_\allowbreak{}evenHeadSynthesis}.\par
\textbf{Recorded historical relation:} \texttt{DIRECT}.\par
\textbf{Current relation:} \texttt{DIRECT}.\par
\textbf{Consumer source state:} \texttt{CURRENT}.\par
\end{lemma}
\begin{lemma}[Historical edge record E004]
\textbf{Supplier:} \texttt{Q3.RouteB.D0Pstar.sourceWeilGraphEvenNonzeroTail\_\allowbreak{}low\_\allowbreak{}fourier\_\allowbreak{}vanish}.\par
\textbf{Consumer:} \texttt{Q3.RouteB.D0Pstar.sourceWeilGraphEvenNonzeroLowBandSynthesis\_\allowbreak{}orthogonal\_\allowbreak{}tail}.\par
\textbf{Recorded historical relation:} \texttt{DIRECT}.\par
\textbf{Current relation:} \texttt{DIRECT}.\par
\textbf{Consumer source state:} \texttt{CURRENT}.\par
\end{lemma}
\begin{lemma}[Historical edge record E005]
\textbf{Supplier:} \texttt{Q3.RouteB.D0Pstar.sourceWeilGraphEvenZeroMode\_\allowbreak{}orthogonal\_\allowbreak{}nonzeroTail}.\par
\textbf{Consumer:} \texttt{Q3.RouteB.D0Pstar.sourceWeilGraphEvenHeadSynthesis\_\allowbreak{}orthogonal\_\allowbreak{}tail}.\par
\textbf{Recorded historical relation:} \texttt{DIRECT}.\par
\textbf{Current relation:} \texttt{DIRECT}.\par
\textbf{Consumer source state:} \texttt{CURRENT}.\par
\end{lemma}
\begin{lemma}[Historical edge record E006]
\textbf{Supplier:} \texttt{Q3.RouteB.D0Pstar.sourceWeilGraphAmbient\_\allowbreak{}evenNonzeroMode\_\allowbreak{}orthonormal}.\par
\textbf{Consumer:} \texttt{Q3.RouteB.D0Pstar.sourceWeilEvenAmbientMode\_\allowbreak{}orthonormal}.\par
\textbf{Recorded historical relation:} \texttt{DIRECT}.\par
\textbf{Current relation:} \texttt{NOT\_\allowbreak{}ESTABLISHED}.\par
\textbf{Consumer source state:} \texttt{HISTORICAL\_\allowbreak{}SOURCE\_\allowbreak{}ADVANCED}.\par
\end{lemma}
\begin{lemma}[Historical edge record E007]
\textbf{Supplier:} \texttt{Q3.RouteB.D0Pstar.sourceWeilGraphAmbient\_\allowbreak{}evenHeadSynthesis}.\par
\textbf{Consumer:} \texttt{Q3.RouteB.D0Pstar.ccmFiniteSynthesis\_\allowbreak{}evenCoefficientEmbedding}.\par
\textbf{Recorded historical relation:} \texttt{DIRECT}.\par
\textbf{Current relation:} \texttt{DIRECT}.\par
\textbf{Consumer source state:} \texttt{CURRENT}.\par
\end{lemma}
\begin{lemma}[Historical edge record E008]
\textbf{Supplier:} \texttt{Q3.RouteB.D0Pstar.sourceWeilGraphEvenHeadSynthesis\_\allowbreak{}orthogonal\_\allowbreak{}tail}.\par
\textbf{Consumer:} \texttt{Q3.RouteB.D0Pstar.ccmFiniteSynthesis\_\allowbreak{}reflectionEven\_\allowbreak{}orthogonal\_\allowbreak{}evenNonzeroTail}.\par
\textbf{Recorded historical relation:} \texttt{DIRECT}.\par
\textbf{Current relation:} \texttt{DIRECT}.\par
\textbf{Consumer source state:} \texttt{CURRENT}.\par
\end{lemma}
\begin{lemma}[Historical edge record E009]
\textbf{Supplier:} \texttt{Q3.RouteB.D0Pstar.sourceWeilGraphTailAmbientCoercive\_\allowbreak{}of\_\allowbreak{}algebraic}.\par
\textbf{Consumer:} \texttt{Q3.RouteB.D0Pstar.sourceWeilEvenTailAmbientCoercive\_\allowbreak{}explicit}.\par
\textbf{Recorded historical relation:} \texttt{DIRECT}.\par
\textbf{Current relation:} \texttt{NOT\_\allowbreak{}ESTABLISHED}.\par
\textbf{Consumer source state:} \texttt{HISTORICAL\_\allowbreak{}SOURCE\_\allowbreak{}ADVANCED}.\par
\end{lemma}

\section{Bibliography}
\nocite{*}
\bibliographystyle{plain}
\bibliography{references}
